\documentclass[11pt]{article}

\usepackage[T1]{fontenc}
\usepackage[utf8]{inputenc}
\usepackage{lmodern}
\usepackage[margin=1.04in]{geometry}
\usepackage{amsmath,amssymb,mathtools,amsthm}
\usepackage{aliascnt}
\usepackage{mathrsfs}
\usepackage{microtype}
\usepackage{enumitem}
\usepackage{xcolor}
\usepackage{booktabs}
\usepackage{url}
\usepackage{hyperref}
\usepackage[capitalize,noabbrev]{cleveref}

\hypersetup{
  colorlinks=true,
  linkcolor=blue!45!black,
  citecolor=green!35!black,
  urlcolor=blue!55!black,
  pdftitle={Relative Exponential Periods on Polynomial Flags: Polynomial Boundary Systems and an Elliptic Comparison Theorem},
  pdfauthor={Christopher D. Long; Antoine-Auguste Le Blanc},
  pdfsubject={Actual boundary systems, elliptic exponential integrals, differential Galois groups, algebraic specialization},
  pdfkeywords={exponential period, polynomial flag, twisted de Rham cohomology, boundary system, elliptic integral, algebraic independence, E-function}
}

\numberwithin{equation}{section}
\setlength{\emergencystretch}{3em}

\newtheorem{theorem}{Theorem}[section]

\newaliascnt{proposition}{theorem}
\newtheorem{proposition}[proposition]{Proposition}
\aliascntresetthe{proposition}

\newaliascnt{lemma}{theorem}
\newtheorem{lemma}[lemma]{Lemma}
\aliascntresetthe{lemma}

\newaliascnt{corollary}{theorem}
\newtheorem{corollary}[corollary]{Corollary}
\aliascntresetthe{corollary}

\theoremstyle{definition}
\newaliascnt{definition}{theorem}
\newtheorem{definition}[definition]{Definition}
\aliascntresetthe{definition}

\newaliascnt{example}{theorem}
\newtheorem{example}[example]{Example}
\aliascntresetthe{example}

\theoremstyle{remark}
\newaliascnt{remark}{theorem}
\newtheorem{remark}[remark]{Remark}
\aliascntresetthe{remark}

\newcommand{\C}{\mathbb C}
\newcommand{\Q}{\mathbb Q}
\newcommand{\Qbar}{\overline{\mathbb Q}}
\newcommand{\R}{\mathbb R}
\newcommand{\Z}{\mathbb Z}
\newcommand{\N}{\mathbb N}
\newcommand{\Aone}{\mathbb A^1}
\newcommand{\Atwo}{\mathbb A^2}
\newcommand{\Gm}{\mathbb G_m}
\newcommand{\Spec}{\operatorname{Spec}}
\newcommand{\Span}{\operatorname{span}}
\newcommand{\trdeg}{\operatorname{trdeg}}
\newcommand{\per}{\operatorname{per}}
\newcommand{\ord}{\operatorname{ord}}
\newcommand{\End}{\operatorname{End}}
\newcommand{\Hom}{\operatorname{Hom}}
\newcommand{\Ext}{\operatorname{Ext}}
\newcommand{\Gal}{\operatorname{Gal}}
\newcommand{\Lie}{\operatorname{Lie}}
\newcommand{\Frac}{\operatorname{Frac}}
\newcommand{\Sym}{\operatorname{Sym}}
\newcommand{\SL}{\operatorname{SL}}
\newcommand{\GL}{\operatorname{GL}}
\newcommand{\Pic}{\operatorname{Pic}}
\newcommand{\Div}{\operatorname{Div}}
\newcommand{\Res}{\operatorname{Res}}
\newcommand{\Tr}{\operatorname{Tr}}
\newcommand{\Crit}{\operatorname{Crit}}
\newcommand{\supp}{\operatorname{supp}}
\newcommand{\rank}{\operatorname{rank}}
\newcommand{\Ecal}{\mathcal E}
\newcommand{\Fcal}{\mathcal F}
\newcommand{\Hcal}{\mathcal H}
\newcommand{\Lcal}{\mathcal L}
\newcommand{\Mcal}{\mathcal M}
\newcommand{\Ocal}{\mathcal O}
\newcommand{\Pcal}{\mathcal P}
\newcommand{\Rcal}{\mathcal R}
\newcommand{\Vcal}{\mathcal V}
\newcommand{\dd}{\,d}
\newcommand{\KZ}{\mathrm{KZ}}

\title{Relative Exponential Periods on Polynomial Flags:\\
Polynomial Boundary Systems and an Elliptic Comparison Theorem}
\author{Christopher D. Long \and Antoine-Auguste Le Blanc\thanks{Le Blanc cryptographic author identifier: \texttt{b6048125\allowbreak{}30b3535b\allowbreak{}23d6dbb0\allowbreak{}f5abe32d\allowbreak{}1ed1a9de\allowbreak{}35fb4f7d\allowbreak{}04d62cff\allowbreak{}08ea19c8}}}
\date{September 2026}


\newcommand{\A}{\mathbb A}
\newcommand{\M}{\mathscr M}
\newcommand{\Ncal}{\mathcal N}
\newcommand{\Tcal}{\mathcal T}
\newcommand{\Hq}{H_q^2}
\newcommand{\Ad}{\operatorname{Ad}}

\begin{document}
\maketitle

\begin{center}
\small
\texttt{galizur@gmail.com}
\end{center}

\begin{abstract}
For a compact polynomially parametrized two-chain $D$ and a polynomial
phase $q$, we construct the differential system of the actual integrals
$\int_D e^{zq}\omega$ using polynomial representatives in twisted algebraic
de Rham cohomology. Its forcing terms are the actual integrals on the fixed
faces of $D$. No spatial denominators or independently normalized punctual
periods are introduced. A separate actual-solution criterion upgrades
nonmembership in the rational linear span of the boundary functions to
algebraic independence over their fraction field, using projective
differential-Galois disjointness and a special-linear affine orbit.

For the triangle with vertices $(-1/2,-1),(1/2,0),(-1/2,1)$ and phase
$q=y^2-x^3+3x$, we prove that the two moments with amplitudes $1$ and $x$
are algebraically independent over the actual face-function field. Their
values are algebraically independent over the corresponding face-value
field at every nonzero algebraic parameter. Every polynomial amplitude
has a unique two-coordinate remainder and an explicit polynomial twisted
Stokes certificate. The geometric nonvanishing is proved on the actual
density by continuation to the two nodal fibers, and the projective
exclusion is proved from the explicitly identified cubic boundary block.
The actual results do not use the compact polynomial-line completeness
theorem; that theorem supplies only the final formal interpretation of
relations among the boundary symbols.

We retain the face-lattice comparison, rooted Picard--Lefschetz theorem,
and denominator-free single-phase comparison on an explicitly observable
locus. For more general collections we give an abstract local
specialization theorem with separate generic-completeness and flatness
hypotheses. Numerical injectivity is stated on the specialized fiber,
not on the unspecialized local functional algebra. General isotypic and
multicolumn geometric comparisons are not asserted without their actual
linear and determinant-line realizations.
\end{abstract}

\medskip
\noindent\textit{2020 Mathematics Subject Classification.}
Primary 11J91; Secondary 11J81, 12H05, 14D07, 14H40, 32S40.

\noindent\textit{Keywords.}
Exponential periods, polynomial flags, twisted de Rham cohomology,
actual boundary systems, elliptic integrals, algebraic independence,
$E$-functions, Kontsevich--Zagier relations.

\setcounter{tocdepth}{1}
\clearpage
\tableofcontents
\clearpage

\section{Introduction and scope}\label{sec:intro}
Put $k=\Qbar$ and $k_{\R}=k\cap\R$. We consider integrals
\[
 I_D(\omega)(z)=\int_D e^{zq}\omega
\]
over compact polynomially parametrized two-chains. The polynomial phase
may be complex in the algebraic construction; real phases are imposed
only where a real one-dimensional phase density is used.

A complete comparison must keep three questions separate. Polynomial
reduction constructs a boundary-forced system. Differential Galois theory
classifies polynomial relations of an actual solution of that system.
Arithmetic lifting transfers a functional classification to algebraic
parameters. None of these questions is answered solely by the list of
composition factors of a holonomic direct image.

\subsection{The polynomial boundary construction}
Over $F_0=k(z)$, put
\[
 d_z=d+z\,dq\wedge,
 \qquad H_q^2=\Omega^2_{F_0[x,y]/F_0}/d_z\Omega^1_{F_0[x,y]/F_0}.
\]
Theorem~\ref{thm:boundary-system} proves that this space has a basis
represented by $z$-independent polynomial forms $\omega_i$ and identities
\[
 q\omega_i=\sum_j A_{ij}(z)\omega_j+d_z\beta_i.
\]
Here $\beta_i$ is polynomial in $x,y$ and rational in $z$. Integrating the
identities on the original chain gives
\[
 U'=AU+b,\qquad
 U_i=\int_D e^{zq}\omega_i,\qquad
 b_i=\int_{\partial D}e^{zq}\beta_i.
\]
The boundary field is generated by all polynomial-amplitude integrals on
these \emph{fixed faces} and their endpoint exponentials. It is not the
field of all compact polynomial-line periods.

Theorem~\ref{thm:criterion} supplies an actual-solution comparison when
the homogeneous group contains $\SL_n$, its adjoint module is excluded
from the boundary tensor category, and the actual column has nonzero
quotient by the finite rational linear boundary space. These are separate
hypotheses; abstract nonzero cohomology alone does not imply the last one.

\subsection{The explicit elliptic theorem}
Let
\[
 q=y^2-x^3+3x,\qquad
 D=\{-1/2\le x\le1/2,\ |y|\le1/2-x\},
 \qquad M_j=\int_D x^j e^{zq}\,dx\,dy.
\]
Set
\begin{align*}
 p(x)&=-x^3+x^2+2x+\tfrac14,&
 P_j(z)&=\int_{-1/2}^{1/2}x^j e^{zp(x)}\,dx,\\
 Q_0(z)&=\int_{-1}^{1}e^{z(y^2-11/8)}\,dy,&
 E_a(z)&=e^{-3z/8},\quad E_b(z)=e^{11z/8}.
\end{align*}
Then the actual face field is
\[
 K_{\partial,0}=k(z)(P_0,P_1,Q_0,E_a,E_b).
\]
Theorems~\ref{thm:elliptic-main} and~\ref{thm:values} prove
\[
 K_{\partial,0}[M_0,M_1]\simeq K_{\partial,0}[X_0,X_1]
\]
and, for every $\xi\in k^\times$,
\[
 K_{\partial,\xi}[M_0(\xi),M_1(\xi)]
 \simeq K_{\partial,\xi}[X_0,X_1],
\]
where $K_{\partial,\xi}$ is the field of the corresponding five boundary
values. Every polynomial amplitude reduces to these two coordinates plus
actual face terms. In particular, its value is algebraic over the
face-value field exactly when its specialized polynomial twisted de Rham
class is zero (Corollary~\ref{cor:value-exactness}).

The proof has an explicit chain. Polynomial identities give the actual
rank-two system. Rational Gelfand--Leray certificates identify its phase
Gauss--Manin matrix. A nonzero nodal variation of the actual density
excludes the rational linear boundary span. The boundary group has only
one nonabelian simple factor, whose adjoint exponential factors differ
from those of the elliptic system. Finally the affine orbit is dense and
the combined $E$-function system is ordinary at every nonzero finite
parameter. No determinant-line generator is needed.

\subsection{Geometry, formal relations, and the general boundary}
The geometric part retains polynomial relative exactness on the observable
locus of Definition~\ref{def:admissible}, using the relative homology
sequence, rooted Picard--Lefschetz variation, and Bonnet's polynomial
relative-exactness theorem. The disjoint-block comparison is stated on
the actual common phase cover where its hypotheses are checked. Rational
Jacobian morphisms and complex linear combinations of them are carefully
distinguished.

The formal interpretation uses the period operations of
Kontsevich--Zagier~\cite{KontsevichZagier} and the compact polynomial-line
comparison of Papers~II--III~\cite{LongLinearKZ,LongPolynomialKZ}; related
exponential-period realizations are discussed in
\cite{CommelinHabeggerHuber}. These inputs concern boundary identities.
The actual independence theorem for the elliptic triangle is independent
of their formal completeness.

For a general collection, Theorem~\ref{thm:abstract-local} isolates the
local arithmetic step. It assumes an explicit generically complete
formal ideal and flatness over a boundary DVR. We also retain the
abstract affine special-linear orbit calculation and a unit-normalized
flat model. Their application to arbitrary flag collections requires
actual linear comparisons and, at full column rank, formally realized
determinant-line maps. No automatic equivalence with a condition on the
original relative monodromy cocycles is asserted.

\part{Geometric functional comparison}
\section{Polynomial flags and phase densities}\label{sec:flags}

\subsection{Polynomial flag chains}

\begin{definition}[Polynomial two-chain]\label{def:poly-chain}
A polynomial two-chain over $k_{\R}$ is a finite $k$-linear combination of
chains
\[
  D=\Phi_*[0,1]^2,
  \qquad
  \Phi\in k_{\R}[u,v]^2.
\]
Its boundary is the alternating sum of the four polynomial face chains.
A polynomial flag chain is a polynomial two-chain for which, after an affine
change of the square coordinates,
\[
  \Phi(u,v)=\bigl(X(u),A(u)+vB(u)\bigr)
\]
with $X,A,B\in k_{\R}[u]$.
\end{definition}

The chain language is preferable to set-theoretic domains: it records
orientation, multiplicity, and algebraic change of variables without
requiring injectivity of $\Phi$.

For a polynomial two-form $\Omega$, define
\[
  F_{D,q,\Omega}(z)=\int_D e^{zq}\Omega.
\]
Pulling back to the square writes this as
\begin{equation}\label{eq:square-form}
  F_{D,q,\Omega}(z)
  =\int_0^1\int_0^1P(u,v)e^{zQ(u,v)}\,du\,dv,
  \qquad P,Q\in k[u,v].
\end{equation}

\subsection{Gelfand--Leray density}

Let $q\in k_{\R}[x,y]$ be nonconstant and let $D$ be a real polynomial flag chain.  Put
$T=q(|D|)\subset\R$, a finite union of compact intervals.  Remove the critical values of $q$, the critical values
of the restrictions of $q$ to the faces, and the phase values of the face
vertices.  The complement is a finite union of open intervals $e$.

For $t\in e$, let
\[
  C_t=q^{-1}(t),
  \qquad
  B_t=C_t\cap|\partial D|,
  \qquad
  \gamma_t=D\cap C_t.
\]
The last object is an oriented relative one-cycle in $(C_t,B_t)$, with the
multiplicities inherited from $D$.

Let $\Omega$ be a polynomial two-form.  On the smooth locus of $q$, choose a
one-form $\alpha$ satisfying
\begin{equation}\label{eq:GL}
  dq\wedge\alpha=\Omega.
\end{equation}
The restriction class of $\alpha$ to $C_t$ is independent of the choice.
Define
\begin{equation}\label{eq:density}
  \rho_e(t)=\int_{\gamma_t}\alpha_t.
\end{equation}

\begin{lemma}[Phase-density formula]\label{lem:density-formula}
After discarding chain components on which the pullback of $\Omega$ is zero,
there is an $L^1$ density $\rho$ on the compact set $T=q(|D|)$ such that
\begin{equation}\label{eq:laplace-density}
  F_{D,q,\Omega}(z)=\int_T e^{zt}\rho(t)\,dt.
\end{equation}
On every regular interval $e$, the restriction $\rho_e$ is real analytic and
is a branch of a period of the algebraic Gauss--Manin system.  In particular,
it admits multivalued holomorphic continuation on the regular phase line.  No
claim that $\rho_e$ is an algebraic function is made.
\end{lemma}

\begin{proof}
On a compact subinterval of $e$, the map $q$ is a submersion on the sliced
chain.  The Gelfand--Leray or coarea formula gives
\[
  \int_{D\cap q^{-1}(e)}e^{zq}\Omega
  =\int_e e^{zt}\left(\int_{\gamma_t}\alpha_t\right)dt.
\]
The inner integral is a period of the relative algebraic de Rham class
$[\alpha_t]$ and therefore is real analytic on $e$ and has the usual
multivalued Gauss--Manin continuation.  The coarea decomposition of the
finite measure obtained by integrating the polynomial density on the compact
chain shows that the resulting density is integrable across the finitely many
critical and boundary values.  Summing the regular intervals proves
\eqref{eq:laplace-density}.  A full-dimensional component on which $q$ is
constant cannot occur for nonconstant polynomial $q$; a lower-dimensional
component pulls every two-form back to zero.
\end{proof}

\subsection{Laplace uniqueness}

\begin{lemma}[Compact-support Laplace uniqueness]\label{lem:laplace-unique}
Let $\rho$ be a compactly supported integrable function on $\R$.  If
\[
  \int_{\R}e^{zt}\rho(t)\,dt=0
  \qquad(z\in\C),
\]
then $\rho=0$ almost everywhere.  If $\rho$ is analytic on each member of a
finite interval decomposition, it vanishes identically on every open
interval.
\end{lemma}

\begin{proof}
Restriction to $z=is$, $s\in\R$, is the Fourier transform of the compactly
supported $L^1$ function $\rho$.  Fourier uniqueness gives almost-everywhere
vanishing, and analyticity gives intervalwise vanishing.
\end{proof}

For multiphase families it is useful to retain the Cauchy-transform form.

\begin{lemma}[Cauchy jump separation]\label{lem:Cauchy-jump}
Let $T\subset\R$ be a finite union of compact intervals and let $\rho_j$ be
piecewise analytic integrable densities on $T$.  If
\[
  \sum_j c_j\int_T e^{zt}\rho_j(t)\,dt=0
  \qquad(z\in\C),
\]
then
\[
  \sum_j c_j\rho_j(t)=0
\]
on every open regular interval.
\end{lemma}

\begin{proof}
All moments of the combined density vanish.  Its Cauchy transform
\[
  \mathcal C(s)=\int_T\frac{\sum_j c_j\rho_j(t)}{s-t}\,dt
\]
therefore vanishes for large $s$ and hence on $\C\setminus T$.  At an interior
point of a regular interval, the two lateral values differ by $2\pi i$ times
the analytic density.  The jump is zero, so the density vanishes.
\end{proof}

\begin{remark}[Why real phase support is imposed]\label{rem:phase-support}
For a general complex-valued polynomial phase on a real two-chain, the
pushforward is a two-dimensional measure in the complex phase plane.  Its
holomorphic moments can all vanish without the measure vanishing.  Thus the
one-variable Cauchy jump argument is unavailable.  The phase-fibered
hypothesis is mathematical content, not a notational convenience.
\end{remark}

\section{Polynomial-flag periods are \texorpdfstring{$E$}{E}-functions}\label{sec:E-functions}

We use Beukers' normalization of Siegel $E$-functions.  An entire function
\[
  F(z)=\sum_{n\ge0}a_n\frac{z^n}{n!}
\]
is an $E$-function if $a_n\in k$, the joint logarithmic height of the first
$n$ coefficients is $O(n)$, and $F$ satisfies a nonzero linear differential
equation over $k(z)$.

\begin{theorem}[Polynomial-flag $E$-function theorem]\label{thm:E-function}
For $P,Q\in k[u,v]$, the function
\[
  F(z)=\int_0^1\int_0^1P(u,v)e^{zQ(u,v)}\,du\,dv
\]
is an $E$-function.  Consequently every polynomial-flag period
$F_{D,q,\Omega}$ is an $E$-function.
\end{theorem}

\begin{proof}
Termwise integration gives
\[
  F(z)=\sum_{n\ge0}a_n\frac{z^n}{n!},
  \qquad
  a_n=\int_0^1\int_0^1P(u,v)Q(u,v)^n\,du\,dv.
\]
Choose a number field $K$ containing the coefficients and an integer $d\ge1$
such that $dP,dQ$ have algebraic-integer coefficients.  The degree of
$PQ^n$ is $O(n)$, the number of its monomials is polynomial in $n$, and every
archimedean conjugate of every coefficient is bounded by $C^n$ for a fixed
$C$.  Writing
\[
  PQ^n=\sum_{i,j} c_{i,j,n}u^iv^j
\]
gives
\[
  a_n=\sum_{i,j}\frac{c_{i,j,n}}{(i+1)(j+1)}.
\]
If $N_n=O(n)$ bounds all occurring exponents, then
\[
  d^{n+1}\operatorname{lcm}(1,\ldots,N_n+1)^2
\]
clears all denominators.  The elementary Chebyshev bound gives
$\log\operatorname{lcm}(1,\ldots,N)=O(N)$; hence the common denominator has
logarithm $O(n)$.  The conjugate and monomial-count estimates give the same
bound for the joint height.

The integrand $P(u,v)e^{zQ(u,v)}$ is holonomic in $(u,v,z)$.  Holonomic
$D$-modules are stable under algebraic direct image, and integration over the
fixed relative chain gives a holonomic function of $z$; see
\cite[Chapters~1 and~3]{HTT}.  Therefore $F$ satisfies a nonzero differential
equation over $k(z)$.  The integral is entire by uniform convergence on
compact subsets.  All three $E$-function conditions hold.
\end{proof}

\begin{theorem}[Beukers lifting]\label{thm:Beukers}
Let $F_1,\ldots,F_N$ be $E$-functions forming a vector solution of a
first-order system over $k(z)$, and let $\xi\in k^\times$ be a nonsingular
point.  Every homogeneous polynomial relation of degree $d$ among the values
$F_i(\xi)$ lifts to a polynomial-coefficient functional relation, homogeneous
of the same degree, specializing to the prescribed relation at $z=\xi$.
\end{theorem}

\begin{proof}
This is the combination of \cite[Theorems~1.3 and~3.2]{Beukers}; affine
relations are homogenized by adjoining the constant $E$-function $1$.
\end{proof}

\section{Relative affine-Galois linearity}\label{sec:relative-AG}

\begin{theorem}[Kolchin--Ostrowski]\label{thm:KO}
Let $K$ be a differential field of characteristic zero with algebraically
closed constants $C$.  Suppose $Y_1,\ldots,Y_m$ lie in a no-new-constants
differential extension and satisfy $DY_i\in K$.  They are algebraically
dependent over $K$ if and only if some nonzero constant linear combination of
them belongs to $K$.
\end{theorem}

This is the classical antiderivative theorem; see
\cite{Kolchin,vanDerPutSinger}.

Let
\[
  R_Y=\{c\in C^m:c\cdot\mathbf Y\in K\}.
\]

\begin{theorem}[Affine-linear presentation]\label{thm:relative-AG}
Let $r=\dim_C R_Y$.  There is $S\in\operatorname{GL}_m(C)$ such that, for
$\mathbf Z=S\mathbf Y$,
\[
  Z_i=g_i\in K
  \quad(1\le i\le r),
\]
and $Z_{r+1},\ldots,Z_m$ are algebraically independent over $K$.  Hence
\[
  \ker\bigl(K[\mathbf X]\to K[\mathbf Y]\bigr)
  =(Z_1-g_1,\ldots,Z_r-g_r).
\]
\end{theorem}

\begin{proof}
Choose a basis of $R_Y$ and extend it to a basis of $C^m$; use these row
vectors to form $S$.  The first $r$ transformed coordinates lie in $K$.  If
the remaining coordinates were algebraically dependent, \cref{thm:KO}
would put a nonzero constant linear combination of them in $K$, contradicting
the choice of complement.  Division by the monic linear polynomials
$Z_i-g_i$ leaves a polynomial in algebraically independent variables, proving
the kernel statement.
\end{proof}

\subsection{One iterated layer}

Let $x$ be an outer variable and let $K$ be a differential field containing
the inner period functions and their $x$-derivatives.  If
\[
  V_i(x,z)=\int_{x_0}^x f_i(t,z)\,dt,
  \qquad f_i\in K,
\]
then \cref{thm:relative-AG} describes the complete polynomial relation ideal
among the partial outer integrals over $K$.

\begin{remark}[Endpoint evaluation is a separate problem]
\label{rem:endpoint-warning}
Take algebraically independent antiderivatives
\[
  V_i(x)=\int_b^x f_i(t)\,dt
\]
over a differential field $K$.  They have no functional algebraic relation
over $K$, but $V_1(b)=\cdots=V_m(b)=0$.  Endpoint evaluation therefore
creates relations not present in the variable-endpoint function field.  In
particular, \cref{thm:relative-AG} cannot by itself justify a value-level
polynomial theorem for completed two-variable periods.
\end{remark}

\section{The face lattice and the generalized Jacobian}\label{sec:face}

Let $\bar f:\bar{\mathcal C}\to S$ be a smooth proper family of connected
curves and let $\pi:B\to S$ be a nonempty reduced divisor finite \'etale over $S$.
Define
\begin{equation}\label{eq:face-lattice}
  \Fcal_B
  =\ker\left(\pi_*\Z_B\xrightarrow{\deg}\Z_S\right).
\end{equation}

\begin{proposition}[Relative homology sequence]\label{prop:relative-homology}
There is a natural exact sequence of integral local systems
\begin{equation}\label{eq:relative-homology-seq}
  0\longrightarrow R_1\bar f_*\Z
  \longrightarrow R_1\bar f_*(\bar{\mathcal C},B;\Z)
  \xrightarrow{\partial}\Fcal_B
  \longrightarrow0.
\end{equation}
\end{proposition}

\begin{proof}
Apply the homology long exact sequence of the pair $(\bar C_s,B_s)$.
Connectedness gives $H_0(\bar C_s,\Z)=\Z$, and the map
$H_0(B_s,\Z)\to H_0(\bar C_s,\Z)$ is the degree map.  The construction is
locally constant in $s$ and therefore yields \eqref{eq:relative-homology-seq}.
\end{proof}

\begin{proposition}[Generalized-Jacobian torus]\label{prop:gen-Jac}
Let $J=\Pic^0(\bar{\mathcal C}/S)$ and
$J_B=\Pic^0(\bar{\mathcal C}/S,B)$.  There is an exact sequence
\begin{equation}\label{eq:gen-Jac-seq}
  1\longrightarrow T_B\longrightarrow J_B\longrightarrow J\longrightarrow0,
  \qquad
  T_B=\Res_{B/S}\Gm/\Gm,
\end{equation}
and
\begin{equation}\label{eq:character-face}
  X^*(T_B)\simeq\Fcal_B.
\end{equation}
\end{proposition}

\begin{proof}
For reduced modulus, the generalized Jacobian is the extension of the ordinary
Jacobian by the quotient of the Weil restriction of $\Gm$ by the diagonal
$\Gm$; see \cite[Chapter~V]{SerreAGCF}.  Characters are contravariant.  The
character sequence associated to
\[
  1\to\Gm\to\Res_{B/S}\Gm\to T_B\to1
\]
is
\[
  0\to X^*(T_B)\to\pi_*\Z_B\xrightarrow{\deg}\Z_S\to0,
\]
which is \eqref{eq:character-face}.
\end{proof}

\begin{remark}[Variance and weights]\label{rem:variance}
The relative-homology quotient $\Fcal_B$ is a lattice of boundary classes.
The same lattice is the character group of the torus $T_B$; the torus itself
is dual.  In mixed-Hodge terms, relative homology and punctured-curve
cohomology carry opposite Tate twists.  The slogan ``the face module is the
toric part'' is correct only after this variance is specified.
\end{remark}

For a polynomial flag slice, $B_t$ is the finite set of intersections with
the oriented faces and
\[
  \partial\gamma_t\in\Fcal_{B,t}.
\]
The equation $\deg\partial\gamma_t=0$ is simultaneously the balanced face
relation and the boundary-of-boundary identity.

\section{Rooted Picard--Lefschetz observability}\label{sec:PL}

Let $H$ be a finite-dimensional vector space and suppose the monodromy group
$G$ is generated by relative Picard--Lefschetz transformations
\begin{equation}\label{eq:relative-PL}
  (T_i-I)\xi
  =\varepsilon_i\langle\widetilde\nu_i,\xi\rangle\nu_i,
  \qquad \varepsilon_i\in\Q^\times.
\end{equation}
Here $\nu_i$ and $\widetilde\nu_i$ are vanishing and dual vanishing cycles.
The nonzero coefficient allows both orientations and positive powers of
local monodromy after a finite cover.
We use the standard relative Picard--Lefschetz formula; see, for example,
\cite[Chapters~2--3]{DimcaSingularities}.
Let $\gamma$ be a relative slice class.

\begin{definition}[Rooted variation graph]\label{def:rooted-graph}
Put
\[
  I_\gamma
  =\{i:\langle\widetilde\nu_i,\gamma\rangle\ne0\}.
\]
Put a directed edge $i\to j$ when
\[
  \langle\widetilde\nu_j,\nu_i\rangle\ne0.
\]
Let $R_\gamma$ be the vertices reachable from $I_\gamma$, and set
\[
  \Vcal_\gamma=\Span\{\nu_i:i\in R_\gamma\},
  \qquad
  \Ocal_\gamma=\Span\{g\gamma-\gamma:g\in G\}.
\]
\end{definition}

\begin{theorem}[Exact rooted orbit theorem]\label{thm:rooted-orbit}
One has
\[
  \Ocal_\gamma=\Vcal_\gamma.
\]
The space is invariant under $G$.
\end{theorem}

\begin{proof}
The space $\Ocal_\gamma$ is invariant because
\[
  h(g\gamma-\gamma)
  =(hg\gamma-\gamma)-(h\gamma-\gamma).
\]
If $i\in I_\gamma$, \eqref{eq:relative-PL} gives $\nu_i\in\Ocal_\gamma$.
If $\nu_i\in\Ocal_\gamma$ and $i\to j$, then
\[
  (T_j-I)\nu_i
  =\varepsilon_j\langle\widetilde\nu_j,\nu_i\rangle\nu_j,
\]
so $\nu_j\in\Ocal_\gamma$.  Hence
$\Vcal_\gamma\subseteq\Ocal_\gamma$.  The same formula shows that every
$T_j$ preserves $\Vcal_\gamma$.

Conversely, for every generator,
\[
  T_j\gamma-\gamma\in\Vcal_\gamma.
\]
If $g\gamma-\gamma\in\Vcal_\gamma$, then
\[
  T_j g\gamma-\gamma
  =T_j(g\gamma-\gamma)+(T_j\gamma-\gamma)
  \in\Vcal_\gamma.
\]
The same argument applies to $T_j^{-1}$, since
$T_j^{-1}-I$ has the same rank-one image.  Induction on word length in the
generators and their inverses gives
$g\gamma-\gamma\in\Vcal_\gamma$ for every $g$, proving the reverse
inclusion.
\end{proof}

On absolute curve homology the intersection graph is undirected.

\begin{corollary}[Connected-graph observability]\label{cor:connected-observable}
Suppose the vanishing cycles span $V$, their intersection graph is connected,
and $I_\gamma\ne\varnothing$.  Then
\[
  \Ocal_\gamma=V.
\]
\end{corollary}

\begin{proof}
Every vertex is reachable from the nonempty active set.  Apply
\cref{thm:rooted-orbit}.
\end{proof}

\begin{proposition}[Absolute irreducibility]\label{prop:absolute-irred}
Assume in addition that $V$ carries a nondegenerate alternating pairing and
that the $T_i$ are the associated symplectic transvections.  Under the
hypotheses of \cref{cor:connected-observable}, the representation on $V$ is
absolutely irreducible.
\end{proposition}

\begin{proof}
Let $K/\Q$ be a characteristic-zero extension and let
$0\ne W\subset V_K$ be invariant.  Choose $0\ne w\in W$.  Since the
vanishing cycles span and the pairing is nondegenerate, some
$\langle w,\nu_i\rangle$ is nonzero.  Then
$(T_i-I)w$ is a nonzero multiple of $\nu_i$, so $\nu_i\in W$.
Connectivity propagates membership to every vanishing cycle.  Thus
$W=V_K$.
\end{proof}

\section{Denominator-free relative exactness}\label{sec:relative-exactness}

\begin{definition}[Admissible observable flag datum]\label{def:admissible}
An admissible observable flag datum is a triple $(D,q,\Omega)$ with $D$ a
polynomial flag chain, $q\in k_{\R}[x,y]$, and
$\Omega\in\Omega^2_{k[x,y]/k}$, satisfying:
\begin{enumerate}[label=\textup{(A\arabic*)},leftmargin=3.2em]
\item $q$ has finite critical locus and every fibre $q^{-1}(t)$ is connected;
\item the generic affine fibre has one place at infinity;
\item after resolving the compactified pencil, every finite singular fibre
      has exactly one ordinary node and the finite critical values are
      distinct; over their complement the fibres form a smooth proper curve
      family, and the face intersections become a finite \'etale divisor after
      a finite base change;
\item the resulting Lefschetz vanishing cycles span $H_1$ of the compactified
      generic fibre and have connected intersection graph; the relative continuation around these nodes obeys the rank-one formula \eqref{eq:relative-PL}, on a common cover labeling the face points;
\item the sliced chain is nonzero on at least one regular phase interval,
      and on every interval on which it is nonzero it has nonzero incidence
      with at least one Lefschetz vanishing cycle.
\end{enumerate}
\end{definition}

The connected-fibre and finite-critical-locus assumptions put $q$ in the
primitive quasi-fibered class of Bonnet.  The one-place-at-infinity condition
removes an additional puncture lattice: it identifies the first homology of
the affine generic fibre with that of its compactification.

\begin{theorem}[Single-phase denominator-free functional theorem]
\label{thm:intro-single}
Let $(D,q,\Omega)$ be an admissible observable flag datum and put
\[
  F(z)=\int_D e^{zq}\Omega.
\]
If $F\equiv0$, then there is a polynomial $a\in k[x,y]$ such that
\[
  \Omega=da\wedge dq.
\]
Consequently
\[
  e^{zq}\Omega=d\bigl(e^{zq}a\,dq\bigr)
\]
and
\[
  F(z)=\int_{\partial D}e^{zq}a\,dq.
\]
The boundary relation is a compact polynomial-line functional relation and is
therefore generated by the pullback--Stokes operations of
\cite{LongLinearKZ}.
\end{theorem}

The algebraic step is
Bonnet's theorem that topological relative exactness equals polynomial
relative exactness for primitive quasi-fibered polynomial maps
\cite{BonnetRelativeExactness}.



\subsection{Bonnet's theorem}

For a polynomial map $q:\C^2\to\C$, a polynomial one-form $\omega$ is
\emph{topologically relatively exact} if its restriction to every generic
fibre has zero integral on every loop.  It is \emph{algebraically relatively
exact} if
\[
  \omega=dR+a\,dq,
  \qquad R,a\in\C[x,y].
\]

\begin{lemma}[Bonnet's hypotheses in relative dimension one]
\label{lem:Bonnet-hypotheses}
Let $q:\C^2\to\C$ be nonconstant.  If its critical locus is finite and every
fibre is connected, then $q$ is primitive and quasi-fibered in Bonnet's
terminology.
\end{lemma}

\begin{proof}
A finite critical locus has codimension two, so $q$ is nonsingular in
codimension one.  Since the target has dimension one, Bonnet's requirement
that the fibres be $1$-generically connected means that every fibre is
connected.  His $2$-generic nonemptiness condition is automatic: for every
$c\in\C$, the proper ideal $(q-c)\subset\C[x,y]$ has a zero by the
Nullstellensatz.  The weaker generic connectedness and nonemptiness conditions
defining primitivity follow as well.
\end{proof}

\begin{theorem}[Bonnet--Gavrilov relative exactness]\label{thm:Bonnet}
If $q:\C^2\to\C$ has finite critical locus and every fibre is connected, then
every topologically relatively exact polynomial one-form is algebraically
relatively exact.
\end{theorem}

\begin{proof}
Apply \cref{lem:Bonnet-hypotheses} and
\cite[Theorem~1.5]{BonnetRelativeExactness}.  Bonnet also records the earlier
two-variable connected-and-reduced theorem of Gavrilov.
\end{proof}

\subsection{From a zero phase density to relative exactness}

Let $\Omega$ be a polynomial two-form.  Choose a polynomial one-form $\omega$
with
\begin{equation}\label{eq:domega}
  d\omega=\Omega;
\end{equation}
this is possible because $H^2_{\mathrm{dR}}(\Atwo)=0$.
On a regular fibre choose $\alpha$ as in \eqref{eq:GL}.  The
Gelfand--Leray formula is
\begin{equation}\label{eq:GM-derivative}
  \nabla[\omega_t]=[\alpha_t].
\end{equation}

\begin{lemma}[Observable density kills the Gelfand--Leray class]
\label{lem:kill-alpha}
Let $(D,q,\Omega)$ satisfy \textup{(A2)--(A5)} of
\cref{def:admissible}.  If the phase density vanishes on a regular interval
on which the sliced chain is nonzero and has nonzero Lefschetz incidence, then
\[
  [\alpha_t]=0\in H^1_{\mathrm{dR}}(C_t)
\]
for the generic affine fibre.
\end{lemma}

\begin{proof}
On the active interval in the hypothesis, the period
$\int_{\gamma_t}\alpha_t$ vanishes and hence, by analytic continuation,
on every continuation of that branch.  Therefore
$[\alpha_t]$ annihilates
\[
  \Span\{g\gamma_t-\gamma_t:g\in G\}.
\]
By \cref{cor:connected-observable}, this is the complete projective
weight-one homology.  The one-place-at-infinity hypothesis identifies this
with $H_1(C_t)$ of the affine fibre.  The de Rham--Betti pairing is perfect,
so $[\alpha_t]=0$.
\end{proof}

\begin{lemma}[The primitive one-form is fibrewise exact]\label{lem:omega-exact}
Under the hypotheses of \cref{lem:kill-alpha}, the restriction of $\omega$ to
the generic fibre is exact.
\end{lemma}

\begin{proof}
Equation \eqref{eq:GM-derivative} and \cref{lem:kill-alpha} show that
$[\omega_t]$ is horizontal.  It is therefore monodromy invariant.  For a
Picard--Lefschetz transvection, invariance of the cohomology class forces its
period on the corresponding vanishing cycle to be zero.  Since the vanishing
cycles span $H_1(C_t)$, all periods of $\omega_t$ vanish.  Hence the
restriction is exact.
\end{proof}

\begin{theorem}[Polynomial relative Stokes]\label{thm:polynomial-relative-Stokes}
Let $(D,q,\Omega)$ be admissible observable.  If the phase density vanishes on
every regular interval, there is $a\in k[x,y]$ such that
\begin{equation}\label{eq:Omega-a}
  \Omega=da\wedge dq.
\end{equation}
Consequently
\begin{equation}\label{eq:twisted-primitive}
  e^{zq}\Omega=d\bigl(e^{zq}a\,dq\bigr).
\end{equation}
\end{theorem}

\begin{proof}
By \textup{(A5)} there is an active regular interval with nonzero Lefschetz
incidence.  The density vanishes there by hypothesis, so
\cref{lem:kill-alpha,lem:omega-exact} show that $\omega$ is topologically
relatively exact.  Apply \cref{thm:Bonnet}:
\[
  \omega=dR+a\,dq
\]
with $R,a\in\C[x,y]$.  Taking exterior derivatives and using
\eqref{eq:domega} gives \eqref{eq:Omega-a}.  The coefficients descend to
$k$: after fixing the finitely many monomials occurring in one solution,
coefficient comparison is a finite linear system over $k$.  Finally,
$dq\wedge dq=0$ gives \eqref{eq:twisted-primitive}.
\end{proof}

\begin{proof}[Proof of \cref{thm:intro-single}]
Functional vanishing and \cref{lem:laplace-unique} make every regular phase
density zero.  Apply \cref{thm:polynomial-relative-Stokes} and ordinary
Stokes:
\[
  F(z)=\int_{\partial D}e^{zq}a\,dq.
\]
Every face is polynomially parametrized, so the right side is a finite
combination of compact polynomial-line functions.  It vanishes identically
because $F\equiv0$.  The functional formal theorem of
\cite{LongLinearKZ} makes this boundary relation formal.
\end{proof}

\begin{remark}[Why one place at infinity appears]\label{rem:infinity}
If the generic affine fibre has $r>1$ points at infinity, then
$H_1(C_t)$ contains an additional rank-$r-1$ puncture lattice.  Projective
Picard--Lefschetz observability does not detect it.  One may replace
\textup{(A2)} by an explicit infinity-lattice observability hypothesis, but it
cannot simply be omitted.
\end{remark}

\section{Hodge rigidity and Jacobian correspondences}\label{sec:Hodge}

\begin{theorem}[Deligne uniqueness for irreducible factors]
\label{thm:Deligne-unique}
An irreducible complex local system underlying a polarizable complex
variation of Hodge structure carries such a variation uniquely up to tensoring
by a rank-one complex Hodge structure $\C(r,s)$, equivalently up to a shift of
the Hodge bigrading.
\end{theorem}

This is \cite[Proposition~1.13]{DeligneFinitude}; see also the convenient
formulation in \cite[Theorem~2]{BakkerMatsushita}.

\begin{theorem}[Weight-one Hodge rigidity]\label{thm:Hodge-rigidity}
Let $\mathbb V$ and $\mathbb W$ be nonzero polarizable rational variations of
Hodge structure of effective weight one on a connected smooth algebraic base.
Assume their complex local systems are irreducible.  Every nonzero flat map
\[
  \phi:\mathbb V_{\C}\longrightarrow\mathbb W_{\C}
\]
is a morphism of complex variations of Hodge structure.
\end{theorem}

\begin{proof}
Irreducibility makes $\phi$ an isomorphism of complex local systems.  Transport
the Hodge bigrading of $\mathbb W$ to $\mathbb V$.  By
\cref{thm:Deligne-unique}, the transported grading differs from the original
one by a shift $\C(r,s)$.  Because both variations are effective of weight
one and arise from rational Hodge structures, both types $(1,0)$ and $(0,1)$
occur.  The shifted set of types can again be exactly
$\{(1,0),(0,1)\}$ only when $r=s=0$.  Thus $\phi$ preserves the Hodge
bigrading and hence the Hodge filtration.
\end{proof}

\begin{remark}[The reducible warning]\label{rem:reducible-Hodge}
The theorem fails for repeated reducible local systems.  A constant elliptic
curve gives a two-dimensional trivial local system; most constant matrices are
flat but do not preserve $H^{1,0}$.  The irreducibility hypothesis is therefore
load-bearing.
\end{remark}

\begin{proposition}[Algebraic realization]\label{prop:Jac-correspondence}
Let $f_i:\mathcal C_i\to S$ and $f_j:\mathcal C_j\to S$ be smooth proper
curve families over a smooth connected complex algebraic base, with relative
Jacobians $J_i,J_j$.  A rational morphism of weight-one variations
\[
  R^1f_{i*}\Q\longrightarrow R^1f_{j*}\Q
\]
is induced, after multiplication by a positive integer and with the usual
contravariance convention, by a homomorphism of abelian schemes between
$J_i$ and $J_j$.  Its action on $R^1$ is represented by an algebraic divisor
correspondence on $\mathcal C_i\times_S\mathcal C_j$.
\end{proposition}

\begin{proof}
The Hodge realization functor from abelian schemes over $S$ to torsion-free
polarizable integral variations of types $(-1,0)$ and $(0,-1)$ is fully
faithful; see \cite[Rappel~(4.4.3)]{DeligneHodgeIII}.  After multiplying the
rational morphism by an integer, full faithfulness therefore supplies the
corresponding homomorphism of relative Jacobians.

The homomorphism and the Poincar\'e bundle determine a relative line bundle on
$J_i\times_S J_j$.  Pulling it back along the Abel--Jacobi maps (after an
inessential finite base change if relative degree-one divisor classes are
needed) gives a line bundle on
$\mathcal C_i\times_S\mathcal C_j$.  Its first Chern class, represented by a
relative divisor with rational coefficients after rescaling, acts on $R^1$ as
the original morphism.
\end{proof}

\begin{corollary}[Complex intertwiners are generated by correspondences]
\label{cor:complex-correspondences}
Under the hypotheses of \cref{thm:Hodge-rigidity}, every complex flat
intertwiner is a complex linear combination of rational VHS morphisms and
therefore, by \cref{prop:Jac-correspondence}, of algebraic Jacobian
correspondences.
\end{corollary}

\begin{proof}
For rational local systems, the intertwiner space is the kernel of a finite
system of rational linear equations.  Hence scalar extension gives
\[
  \Hom_{\pi_1}(\mathbb V,\mathbb W)\otimes_{\Q}\C
  \simeq
  \Hom_{\pi_1}(\mathbb V_{\C},\mathbb W_{\C}).
\]
Apply \cref{thm:Hodge-rigidity,prop:Jac-correspondence} to the rational
summands.
\end{proof}

\section{Multiphase functional comparison}\label{sec:multiphase}

Let $S$ be the phase line with all critical and boundary branch values removed.
After a common finite \'etale base change, assume that all face points are
sections.  For each datum let
\[
  0\to V_j\to H_j^{\mathrm{rel}}\to\Fcal_j\to0
\]
be the relative-homology sequence and let $\gamma_j$ be the sliced relative
class.  Let $\alpha_j$ denote the Gelfand--Leray cohomology class.

\subsection{The direct-sum variation space}

Set
\[
  V=\bigoplus_{j=1}^N V_j,
  \qquad
  \Gamma=(\gamma_1,\ldots,\gamma_N),
\]
and
\begin{equation}\label{eq:combined-orbit}
  \Ocal_\Gamma
  =\Span\{g\Gamma-\Gamma:g\in\pi_1(S)\}\subseteq V.
\end{equation}
The projection of $\Ocal_\Gamma$ to $V_j$ is $\Ocal_{\gamma_j}$ and hence,
for an observable datum, is all of $V_j$.

\begin{lemma}[Disjoint irreducible factors]\label{lem:disjoint-factors}
Let $V_1,\ldots,V_N$ be pairwise nonisomorphic irreducible complex
representations of a group $G$.  If a $G$-submodule
$W\subseteq\bigoplus_j V_j$ projects nontrivially to every $V_j$, then
\[
  W=\bigoplus_j V_j.
\]
\end{lemma}

\begin{proof}
The finite-dimensional image algebra of $\C[G]$ is semisimple on the direct
sum.  A submodule is therefore a direct sum of some of the pairwise
nonisomorphic simple factors.  Nonzero projection to each factor forces every
factor to occur.
\end{proof}

\subsection{Proof in the disjoint-block case}

\begin{theorem}[Disjoint-block functional comparison]
\label{thm:disjoint-comparison}
Let $(D_j,q_j,\Omega_j)$ be admissible observable data. On the chosen
common regular phase cover labeling their face points, suppose that the
systems $V_{j,\C}$ are irreducible, pairwise nonisomorphic, and that every
active slice remains observable. Every relation
\[
 \sum_j c_j\int_{D_j} e^{zq_j}\Omega_j\equiv0,\qquad c_j\in k,
\]
is generated by linearity, polynomial relative Stokes, and compact
polynomial-line functional relations on the faces. The hypotheses concern
the representations on this common cover, not only on an earlier base.
\end{theorem}

\begin{proof}
By \cref{lem:Cauchy-jump}, the combined phase density vanishes on every open
regular interval.  Fix such an interval $e$ and let $J_e$ be the indices for
which the sliced chain is nonzero on $e$.  Analytic continuation to a base
point of the common regular phase line gives, for every $g\in\pi_1(S)$,
\[
  \sum_{j\in J_e}c_j\langle\alpha_j,g\gamma_{j,e}\rangle=0.
\]
Subtract the equation for $g=1$.  The covector
\[
  \alpha_e=(c_j\alpha_j)_{j\in J_e}
  \in\left(\bigoplus_{j\in J_e}V_j\right)^\vee
\]
annihilates the orbit-variation submodule associated with the tuple
$(\gamma_{j,e})_{j\in J_e}$.  By admissibility, its projection to every
$V_j$, $j\in J_e$, is surjective.  The factors remain pairwise
nonisomorphic, so \cref{lem:disjoint-factors} makes the orbit submodule the
whole direct sum.  Hence $c_j\alpha_j=0$ for every $j\in J_e$.

Every nonzero two-chain term is active on at least one regular interval.
Repeating the argument over the finite interval decomposition therefore gives
$c_j\alpha_j=0$ for every index with $c_j\ne0$.  The proof of
\cref{thm:polynomial-relative-Stokes} then applies to the scaled form
$c_j\Omega_j$ and gives a polynomial $a_j$ with
\[
  c_j\Omega_j=da_j\wedge dq_j.
\]
Stokes reduces the relation to the faces.  The face terms are compact
polynomial-line functional relations; they are generated by the formal
operations of \cite{LongLinearKZ}.  Their zero-dimensional endpoint terms
are already part of the compact polynomial-line formal space.
\end{proof}


\begin{remark}[What a Jacobian correspondence does not supply]
A rational weight-one morphism has the algebraic realization above, but a
complex scalar multiple need not become rational after multiplication by an
integer. For example, $\sqrt2\,\mathrm{id}$ never does. Moreover, an
isomorphism of generic absolute systems does not identify actual relative
slice classes or preserve a polynomial lattice in their de Rham
representatives. The general isotypic formal comparison is not asserted
here without those additional comparisons. None of them is assumed in the
elliptic theorem below.
\end{remark}

\part{Actual boundary systems and elliptic comparison}
\section{Polynomial representatives and literal boundary forcing}
Put $k=\overline{\Q}$ and $F_0=k(z)$. Let $D$ be a finite algebraic-linear
combination of pushforwards of $[0,1]^2$ by polynomial maps with real
algebraic coefficients. Let $q\in k[x,y]$. For a polynomial two-form
$\omega$, write
\[
 I_D(\omega)=\int_D e^{zq}\omega.
\]
When coefficients of $\omega$ lie in $F_0$, interpret this as a meromorphic
function of $z$ by linearity. For each oriented face parametrization
$\phi_e:[a_e,b_e]\to\A^2$, put $p_e=q\circ\phi_e$.
The actual boundary field $K_{\partial,0}$ is the differential field over
$F_0$ generated by all functions
\[
 \int_{a_e}^{b_e} h(t)e^{zp_e(t)}\dd t,\qquad h\in k[t],
\]
and the exponential functions at their algebraic endpoints. Only the fixed
faces are used; this does not mean the field of all compact line periods.

Let $d$ differentiate only in $x,y$, and set
\[
 d_z=d+z\,dq\wedge,\qquad
 \Hq=\Omega^2_{F_0[x,y]/F_0}/d_z\Omega^1_{F_0[x,y]/F_0}.
\]

\begin{theorem}[Polynomial-representative boundary system]
\label{thm:boundary-system}
The space $\Hq$ is finite dimensional. There are forms
$\omega_i=r_i(x,y)\,dx\wedge dy$, with $r_i\in k[x,y]$, whose classes form
an $F_0$-basis, a matrix $A\in M_n(F_0)$, and polynomial-in-$x,y$ one-forms
$\beta_i\in\Omega^1_{F_0[x,y]/F_0}$ such that
\begin{equation}\label{eq:reduction-system}
 q\omega_i=\sum_jA_{ij}(z)\omega_j+d_z\beta_i.
\end{equation}
The actual integral column $U_i=I_D(\omega_i)$ satisfies
\begin{equation}\label{eq:actual-system}
 U'=AU+b,\qquad b_i=\int_{\partial D}e^{zq}\beta_i\in K_{\partial,0}.
\end{equation}
Every polynomial two-form $\omega$ has an identity
\begin{equation}\label{eq:amplitude-normal}
 \omega=\sum_i a_i(z)\omega_i+d_z\beta,
\end{equation}
where $a_i\in F_0$ and $\beta$ is polynomial in $x,y$. Consequently
\begin{equation}\label{eq:actual-normal}
 I_D(\omega)=\sum_i a_i(z)U_i+
              \int_{\partial D}e^{zq}\beta.
\end{equation}
For any finite set of amplitudes, these functions belong to a finite closed
rational first-order system whose coordinates are actual $E$-functions.
No spatial denominator is introduced. Parameter denominators may occur.
\end{theorem}

\begin{proof}
Algebraic direct image of the rank-one connection $d+ d(zq)$ on
$\A^2\times\A^1_z$ is holonomic. The generic top relative de Rham
cohomology is therefore finite dimensional over $k(z)$. Equivalently, it is
the generic localization of the Fourier transform of the polynomial
Gauss--Manin system. The polynomial complex and Fourier convention are
spelled out in \cite[Proposition 9.1 and Remarks 9.4]{Sabbah}, with Fourier
parameter $\tau=-z$. This finiteness statement does not identify the rank
with a Milnor number without an additional hypothesis.

The monomial two-forms with coefficients in $k$ span the quotient over
$F_0$, so a finite subset is a basis. The connection on the quotient is
\[
 \nabla_z[\omega]=[\partial_z\omega+q\omega],
\]
because
\[
 [\partial_z+q,d+z\,dq\wedge]
 =dq\wedge-dq\wedge=0.
\]
Expressing $q\omega_i$ in the basis proves \eqref{eq:reduction-system};
expressing an arbitrary class proves \eqref{eq:amplitude-normal}.

For a one-form polynomial in the spatial variables, there is the literal
identity
\[
 e^{zq}d_z\beta=d_{x,y}(e^{zq}\beta).
\]
Stokes on the original parametrized compact chain proves
\eqref{eq:actual-system} and \eqref{eq:actual-normal}. Each pullback of
$\beta$ to a face is a finite $F_0$-linear combination of polynomial
one-forms in its parameter. Thus every forcing term is an actual face
integral with coefficients in $F_0$.

To see finiteness on the boundary, if $\deg p=d\ge2$, division by
$\partial_t+zp'(t)$ reduces every polynomial amplitude to degree at most
$d-2$. The leading degree of $(\partial_t+zp')h$, for $h\ne0$, is
$\deg h+d-1$, with nonzero leading coefficient in $F_0$. The remainder is
therefore unique. Integration of the exact part gives only endpoint
exponentials. Multiplication by $p$ and the same division give a rational
first-order system for the reduced moments and the endpoint exponentials.
For $d\le1$ every polynomial amplitude reduces directly to endpoint terms.
For finitely many faces, this supplies a finite boundary vector; its
fraction field is $K_{\partial,0}$.

The integral coordinates are entire and have algebraic Taylor coefficients.
After pullback to the square they have the form
$\int_{[0,1]^2}P(u,v)e^{zQ(u,v)}\dd u\dd v$. For the first $m$ Taylor
coefficients, one denominator of the form
\[
 d^{m+1}\operatorname{lcm}(1,\ldots,Cm+C_0)^2
\]
suffices, with fixed integers $d,C,C_0$; all conjugates grow exponentially.
The same argument in one variable handles the faces. Together with the
constructed differential systems, these are the $E$-function conditions.
\end{proof}

\begin{proposition}[Actual affine comparison under a change of representatives]
\label{prop:gauge}
Suppose
$\widetilde\omega=T(z)\omega+d_z\gamma$, with $T\in\GL_n(F_0)$.
Put $g=\int_{\partial D}e^{zq}\gamma$. Then the actual columns satisfy
\[
 \widetilde U=TU+g,
\]
and their systems are related by
\[
 \widetilde A=(T'+TA)T^{-1},\qquad
 \widetilde b=Tb+g'-\widetilde A g.
\]
In particular, this is a comparison of the actual affine solution objects,
not a statement only about their homogeneous composition factors.
\end{proposition}
\begin{proof}
The first equality is Stokes, and differentiating it proves the other two.
If two primitives give the same form identity, their difference is
$d_z$-closed and its integral over $\partial D$ is zero by the same Stokes
identity. Thus primitive choices do not introduce an arbitrary constant.
\end{proof}

After adjoining complex constants, write $F=\C(z)$,
$K_\partial=\C K_{\partial,0}$, and work in a field $\M$ of meromorphic
germs containing the actual functions. Integration modulo the actual
boundary field gives a canonical differential map
\begin{equation}\label{eq:canonical-comparison}
 K_\partial\otimes_{F_0}\Hq\longrightarrow\M/K_\partial,
 \qquad [\omega]\longmapsto I_D(\omega)\bmod K_\partial.
\end{equation}
Here $\M/K_\partial$ is used as a differential vector space, not as a
quotient field or ring. This map need not be injective for arbitrary chains.
The next criterion proves injectivity rather than building it into a
completion definition.

\section{An actual-solution comparison criterion}
Let $\mathcal B$ be a finite rational boundary differential system, let
$L_\partial/F$ be a Picard--Vessiot field for that system containing its
actual solution coordinates, and let $\Ncal\subset L_\partial$ be their
finite-dimensional $F$-linear differential span. Include $1$ when needed.
Let $\Tcal_\partial$ be the tensor category generated by $\mathcal B$,
allowing subquotients and duals. Its objects are precisely the differential
modules trivialized by $L_\partial$ \cite[Theorem 1.5.2]{Singer}.

\begin{lemma}[Differentially finite elements of a Picard--Vessiot field]
\label{lem:dfinite}
Every finite-dimensional $F$-differential subspace of $L_\partial$ belongs
to $\Tcal_\partial$.
\end{lemma}
\begin{proof}
Let $y\in L_\partial$ satisfy a scalar differential equation over $F$.
Every Galois conjugate of $y$ satisfies the same equation. Since there are
no new constants, their constant span $V$ is finite dimensional. Choose a
basis $y_1,\ldots,y_r$ of $V$. Its Wronskian is nonzero, and
\[
  \mathcal L_V(f)=
  \frac{\operatorname{Wr}(y_1,\ldots,y_r,f)}
       {\operatorname{Wr}(y_1,\ldots,y_r)}
\]
is a monic scalar operator of order $r$. Its coefficients lie in
$L_\partial$ and are fixed by its entire Galois group: a basis change of
$V$ multiplies numerator and denominator by the same nonzero constant.
Thus $\mathcal L_V\in F[\partial_z]$. It has a full fundamental set in
$L_\partial$, so its differential module is trivialized there. The
$F$-differential span of $y$ is a quotient of its coefficient module and
belongs to $\Tcal_\partial$. Applying this construction to finitely many
basis elements proves the assertion, since the category is closed under
finite direct sums and quotients. See \cite[Theorems~1.3.9 and~1.5.2]{Singer}
for the fixed-field and tensor correspondences.
\end{proof}

\begin{theorem}[Actual one-column comparison]
\label{thm:criterion}
Suppose an actual column $U$ satisfies $U'=AU+b$, with $b_i\in\Ncal$,
and $n\ge2$. Assume:
\begin{enumerate}
\item the homogeneous Galois group over $F$ contains $\SL_n$;
\item its adjoint module $\End^0(A)$ does not belong to $\Tcal_\partial$;
\item at least one coordinate of $U$ is not in $\Ncal$.
\end{enumerate}
Then $U_1,\ldots,U_n$ are algebraically independent over
$K_\partial\subset L_\partial$. In the setting of
\cref{thm:boundary-system}, \eqref{eq:canonical-comparison} is injective,
and its image is the span of the actual polynomial-amplitude periods
modulo $K_\partial$.
\end{theorem}

\begin{proof}
Let $E=\Ncal+\sum_i FU_i$. The differential module $E/\Ncal$ is a
quotient of the coefficient module associated with $A$. This module is
irreducible because the homogeneous group contains $\SL_n$; taking a dual
according to the convention for connection matrices does not affect the
argument. By the third assumption, the quotient is nonzero and hence is
that entire irreducible module.

If all $U_i$ belonged to $L_\partial$, \cref{lem:dfinite} would put $E$,
then $E/\Ncal$, and then its adjoint module in $\Tcal_\partial$,
contradicting the second assumption. Thus $U\notin L_\partial^n$.

Choose the joint Picard--Vessiot field $L_hL_\partial/F$ inside a
common no-new-constants differential overfield. Restriction onto the
homogeneous group $G$ is surjective. The kernel of restriction to
$L_\partial$ therefore maps to a closed normal subgroup $H\lhd G$.
It is the homogeneous group after base change to $L_\partial$.

Here is the needed normal-subgroup dichotomy, including finite components.
Since $G\subset\GL_n$ contains $\SL_n$, its derived subgroup is $\SL_n$.
The normal subgroup $H\cap\SL_n$ is either all of $\SL_n$ or is finite
central. In the latter case, for fixed $h\in H$ the commutator morphism
$\SL_n\to H\cap\SL_n$, $s\mapsto[h,s]$, is constant by connectedness
and is the identity at $s=I$. Thus $h$ centralizes $\SL_n$ and is scalar.
Hence either $H\supset\SL_n$ or $H$ acts trivially on $\End^0(A)$.
The latter would trivialize the adjoint module over $L_\partial$, contrary
to the second assumption. It follows that the homogeneous group over
$L_\partial$, and hence over $K_\partial$, contains $\SL_n$.
Every subgroup of $\GL_n$ containing $\SL_n$ is reductive: its unipotent
radical intersects $\SL_n$ trivially and has trivial image in the
multiplicative determinant quotient, hence is trivial.

Let $Y$ be a fundamental matrix over a joint field for the inhomogeneous
system over $K_\partial$. Its affine Galois group is a subgroup of
$\C^n\rtimes G_{K_\partial}$, acting by
$\sigma(U)=U+Yc_\sigma$. Its translation kernel is an invariant vector
subspace, hence either $0$ or $\C^n$.
If it is $\C^n$, translations make the orbit of $U$ Zariski dense in
$\A^n$. If it is zero, the remaining regular algebraic cocycle of the
reductive homogeneous group is a coboundary. Thus
\[
 U=a+Yc,\qquad a\in K_\partial^n,\quad c\in\C^n.
\]
The vector $c$ is nonzero, since otherwise $U\in K_\partial^n$. The
$\SL_n$-orbit of a nonzero vector is $\A^n\setminus\{0\}$, again dense.
In either case a polynomial over $K_\partial$ vanishing at $U$ vanishes
on a dense affine orbit and is zero. Linear independence modulo the
boundary field, and then the last assertion by normal-form reduction,
follow immediately.
\end{proof}

\begin{remark}
The third condition concerns only rational \emph{linear} combinations of a
finite actual boundary vector. It is strictly weaker than the desired
nonmembership in the boundary fraction field. The proof is what upgrades
this linear test to the fraction-field statement. A nonzero abstract
cohomology class alone would not suffice.
\end{remark}

\section{The elliptic polynomial normal form}
Fix the counterclockwise triangle
\[
 D=\operatorname{conv}\{(-1/2,-1),(1/2,0),(-1/2,1)\}
 =\{-1/2\le x\le1/2,\ |y|\le1/2-x\},
\]
and
\[
 q=y^2-x^3+3x,\qquad
 M_j=\int_Dx^j e^{zq}\dd x\dd y\quad(j=0,1).
\]
A polynomial parametrization is
$\Phi(u,v)=(u-1/2,(1-u)(2v-1))$, with Jacobian $2(1-u)$.

For $\beta=P\,dy-Q\,dx$, one has
\[
 d_z\beta=(L_xP+L_yQ)\,dx\wedge dy,
\quad L_x=\partial_x+3z(1-x^2),\quad L_y=\partial_y+2zy.
\]
The operators act in separate variables. Their one-variable quotients have
bases $1,x$ and $1$, respectively; both operators are injective on
polynomials. Tensoring the two one-variable complexes proves
\begin{equation}\label{eq:ranktwo}
 \Hq=F_0[dx\wedge dy]\oplus F_0[x\,dx\wedge dy].
\end{equation}
In particular the rank is exactly two, not merely at most two.

The explicit reduction recurrences are
\begin{align}
 [x^i y^j]&=-\frac{j-1}{2z}[x^i y^{j-2}] &&(j\ge1),\label{eq:yrec}\\
 [x^{m+2}]&=[x^m]+\frac{m}{3z}[x^{m-1}] &&(m\ge0),\label{eq:xrec}
\end{align}
with the zero-coefficient terms omitted. Each use records a polynomial
primitive. Hence every $r\in k[x,y]$ has a unique remainder
$a_r(z)+b_r(z)x$, with $a_r,b_r\in k[z,z^{-1}]$, and an identity
\begin{equation}\label{eq:elliptic-normal}
 r=a_r+b_rx+L_xP_r+L_yQ_r,
 \qquad P_r,Q_r\in k[z,z^{-1},x,y].
\end{equation}

Let
\[
 C=\begin{pmatrix}0&2\\2&0\end{pmatrix},\qquad
 D_0=\begin{pmatrix}-5/6&0\\0&-7/6\end{pmatrix},\qquad A=C+D_0/z.
\]
The following are identities of polynomial forms, with Laurent coefficients
in $z$:
\begin{align}
 q\,dx\wedge dy
 &= -\frac5{6z}\,dx\wedge dy+2x\,dx\wedge dy+d_z\beta_0,\\
 qx\,dx\wedge dy
 &=2\,dx\wedge dy-\frac7{6z}x\,dx\wedge dy+d_z\beta_1,\label{eq:explicit-beta}
\end{align}
where
\[
 \beta_0=\frac{x}{3z}\,dy-\frac{y}{2z}\,dx,
 \qquad
 \beta_1=\frac{x^2-2}{3z}\,dy-\frac{xy}{2z}\,dx.
\]
Define actual face integrals
\[
 B_{x,j}=\int_{\partial D}x^je^{zq}\,dy,
 \qquad
 B_{y,j}=-\int_{\partial D}x^jye^{zq}\,dx.
\]
Then
\begin{equation}\label{eq:elliptic-system}
 \binom{M_0}{M_1}'
 =\begin{pmatrix}-5/(6z)&2\\2&-7/(6z)\end{pmatrix}\binom{M_0}{M_1}
 +\frac1z\binom{B_{y,0}/2+B_{x,1}/3}
 {B_{y,1}/2+B_{x,2}/3-2B_{x,0}/3}.
\end{equation}
This is obtained by integrating \eqref{eq:explicit-beta}, not by choosing
an isomorphic abstract connection and then normalizing a solution.

\subsection{A finite actual boundary field}
Put
\[
 p(x)=-x^3+x^2+2x+\tfrac14,
 \quad P_j(z)=\int_{-1/2}^{1/2}x^j e^{zp(x)}\dd x,
 \quad Q_0(z)=\int_{-1}^{1}e^{z(y^2-11/8)}\dd y,
\]
and $E_a=e^{-3z/8}$, $E_b=e^{11z/8}$. Parametrizing the three oriented
edges gives
\begin{equation}\label{eq:face-identities}
 B_{x,j}=2P_j-(-1/2)^jQ_0,
 \qquad B_{y,j}=P_j-2P_{j+1}.
\end{equation}
The one-variable reductions imply
\begin{equation}\label{eq:boundary-field}
 K_{\partial,0}=k(z)(P_0,P_1,Q_0,E_a,E_b).
\end{equation}
Indeed the two slanted faces have the same phase $p$; their arbitrary
polynomial amplitudes reduce to $P_0,P_1$ and endpoints. The vertical
quadratic amplitudes reduce to $Q_0$ and endpoints.

For completeness, a closed boundary system is
\begin{align}
 P_0'&=\left(\frac{17}{36}-\frac1{3z}\right)P_0+
       \frac{14}{9}P_1+\frac{E_b+5E_a}{18z},\\
 P_1'&=\left(\frac{28}{27}+\frac1{9z}\right)P_0+
       \left(\frac{163}{108}-\frac2{3z}\right)P_1+
       \frac{-53E_b+41E_a}{108z},\\
 Q_0'&=\left(-\frac{11}{8}-\frac1{2z}\right)Q_0+\frac{E_a}{z},
 \qquad E_a'=-\tfrac38 E_a,\quad E_b'=\tfrac{11}{8}E_b.
\end{align}
For example, $P_2=(2P_1+2P_0)/3-(E_b-E_a)/(3z)$, so
\eqref{eq:elliptic-system} closes with this boundary vector. The only
finite pole of the combined rational system is zero.

\subsection{The actual normalization at zero}
The face forcing in \eqref{eq:elliptic-system} is $b(z)/z$ with $b$ entire.
Writing $M=\sum v_mz^m$, $b=\sum b_mz^m$, its holomorphic solution obeys
\[
 (mI-D_0)v_m=Cv_{m-1}+b_m,\qquad v_{-1}=0.
\]
Every $mI-D_0$, $m\ge0$, is invertible. Thus there is at most one
holomorphic solution at zero with these actual forcing functions. The
integral supplies it, with
\[
 M_0(0)=1,\qquad M_1(0)=-\tfrac16.
\]
This supplies a second, independent check that no unspecified homogeneous
solution has been added to the actual column.

\section{Explicit comparison with elliptic Gauss--Manin periods}
The regular fibers are
\[
 E_t:\ y^2=f_t(x)=x^3-3x+t.
\]
On a locally constant closed cycle, let
$h_j(t)=\int x^j dx/(2y)$ for $j=0,1$. Set
\[
 R_0(x,t)=\frac{4-tx-2x^2}{6(t^2-4)},\qquad
 R_1(x,t)=\frac{2t-2x-tx^2}{6(t^2-4)}.
\]
Direct differentiation of rational forms gives
\begin{equation}\label{eq:GLcertificate}
 \partial_t\left(\frac{x^jdx}{2y}\right)
 -\sum_{k=0}^1G_{jk}(t)\frac{x^kdx}{2y}
 =d_x\left(\frac{R_j(x,t)}{y}\right),
 \qquad
 G(t)=\frac1{6(t^2-4)}\begin{pmatrix}-t&2\\-2&t\end{pmatrix}.
\end{equation}
These identities can be checked after multiplication by $4y^3$.
Integration on closed cycles proves
\begin{equation}\label{eq:phase-GM}
 h'=G(t)h,\qquad (tI-C)h'=-(D_0+I)h.
\end{equation}
The forms $dx/(2y),x\,dx/(2y)$ are the usual two-dimensional de Rham
basis of the elliptic curve (the latter is of the second kind at infinity).
The exact rational-form identities identify the actual Gauss--Manin
connection, not only its local exponents.

With Laplace convention
$\mathcal L(tT)=\partial_z\mathcal L(T)$ and
$\mathcal L(\partial_tT)=-z\mathcal L(T)$, applying the Weyl substitution
to \eqref{eq:phase-GM} gives
\[
 -\partial_z(zU)-C(-zU)=-(D_0+I)U,
 \quad\text{hence}\quad zU'=zCU+D_0U.
\]
This is exactly the homogeneous part of \eqref{eq:elliptic-system}.
The polynomial certificates \eqref{eq:explicit-beta} give its actual
relative forcing. These two calculations identify the displayed Weyl presentation and the
actual polynomial relative forcing. A connection on the punctured phase
line alone would not specify its extension across the nodes; no uniqueness
of an unspecified holonomic completion is being asserted or used. The
independence argument below uses the explicit matrix \eqref{eq:elliptic-system}
and the actual density, not such a uniqueness claim.

\section{The actual column is not linearly generated by the faces}
In this section all linear coefficients may be complex. Let
\[
 \Ncal=\Span_{\C(z)}\{1,P_0,P_1,Q_0,E_a,E_b\}.
\]
It is a differential space by the displayed boundary system.

\begin{lemma}[Nonzero actual quotient]\label{lem:nonzero-quotient}
One has $M_0\notin\Ncal$.
\end{lemma}
\begin{proof}
Suppose otherwise. Clearing rational denominators yields
\[
 P(z)M_0(z)=\sum_jR_j(z)g_j(z),\qquad
 0\ne P\in\C[z],\quad R_j\in\C[z],\quad
 g_j\in\{1,P_0,P_1,Q_0,E_a,E_b\}.
\]
The inverse Laplace distributions of the actual line periods have
algebraic densities on every regular real interval: on a monotone branch
of a polynomial phase the density is $h(x(t))/p'(x(t))$.
The exponentials and $1$ give point masses. Derivatives of these
distributions still have algebraic germs on regular intervals.
Laplace uniqueness for compactly supported distributions therefore implies
that the density germ $\rho(t)$ of $M_0$, on a regular interval just to
the right of $t=0$, satisfies
\begin{equation}\label{eq:rho-algebraic}
 P(-\partial_t)\rho(t)\quad\text{is algebraic over }\C(t).
\end{equation}

We verify a nonzero elliptic monodromy variation of this particular germ.
Put $f(x)=-x^3+3x$, let $r_2(s)\in(-1,1)$ be its middle inverse branch
for $-2<s<2$, and let $c(t)$ be the inverse branch of $p$ beginning at
$c(0)\in(-1/5,0)$. Set $a(t)=1/2-c(t)$. Coarea on the triangle gives,
near zero,
\begin{equation}\label{eq:actual-density}
 \rho(t)=\int_{-a(t)}^{a(t)}
 \frac{dy}{3(1-r_2(t-y^2)^2)}.
\end{equation}
We check the continuation and its normalization explicitly. On
$[-1/5,1]$ one has $p'>0$, with
\[
 p(-1/5)=-51/500<0<p(0)=1/4,\qquad
 p(4/5)=989/500<2<p(1)=9/4.
\]
Thus the branch $c(t)$ is analytic on $0\le t\le2$,
$c(0)\in(-1/5,0)$, and $c(2)\in(4/5,1)$. For $0<t<2$ and
$|y|\le|a(t)|$,
\[
  t-y^2\in[f(c(t)),t]\subset(-2,2),
  \qquad f(c(t))=t-a(t)^2.
\]
Consequently the inverse branch $r_2$ in \eqref{eq:actual-density} is
unambiguous throughout this real continuation. At $t_*=11/8$ the
endpoints cross because $c(t_*)=1/2$. The expression
\[
 a(t)\int_{-1}^1
 \frac{dv}{3(1-r_2(t-a(t)^2v^2)^2)}
\]
is holomorphic near $t_*$, proving that the crossing does not obstruct
analytic continuation. For $t>t_*$ its orientation is reversed.

At $t=2$ the endpoints remain distinct and unramified, since
$a(2)<-3/10$ and $p'(c(2))\ne0$. Put $\varepsilon=2-t$ and
$r_2(t-y^2)=1-v$. The identity
\[
 f(1-v)=2-3v^2+v^3
\]
gives, uniformly near the node on the chosen branch,
\[
 \frac1{3(1-r_2(t-y^2)^2)}
 =\frac1{2\sqrt3\sqrt{\varepsilon+y^2}}+O(1).
\]
Splitting the integral at a small fixed neighborhood of $y=0$, and using
the reversed orientation, yields the nonzero coefficient
\begin{equation}\label{eq:log-coefficient}
 \rho(2-\varepsilon)
 =\frac1{2\sqrt3}\log\varepsilon+O(1),
 \qquad \varepsilon\downarrow0.
\end{equation}
Outside that neighborhood the relative path and its endpoints transport
holomorphically around a small loop about $2$. Inside it the nodal neck is
an annulus. Its monodromy changes a crossing relative path by a multiple
of its core circle; this is the local Picard--Lefschetz description
\cite[Chapters~2--3]{DimcaSingularities}. The coefficient cannot vanish by
\eqref{eq:log-coefficient}. With $\delta_+$ the core vanishing cycle,
\[
 (T_+-I)\rho=c_*h_+(t),\qquad c_*\ne0,
 \quad h_+(t)=\int_{\delta_+}\frac{dx}{2y},\quad T_+h_+=h_+.
\]
The sign from $dx/(2y)=-dy/(3(1-x^2))$ on a fiber is absorbed in $c_*$.
This fixes the comparison with the actual oriented density and does not
choose an independent normalization of a punctual distribution.

A power $T_+^m$ fixes the algebraic germ in \eqref{eq:rho-algebraic}.
Analytic continuation commutes with the constant-coefficient operator,
so $P(-\partial_t)h_+=0$.
But every solution germ of a nonzero constant-coefficient differential
equation extends to an entire function of $t$.
The elliptic period $h_+$ does not. Indeed, for $-2<t<2$ let
$r_1<r_2<r_3$ be the roots of $x^3-3x+t$. Up to a fixed nonzero factor,
\[
 h_+(t)=\int_{r_2}^{r_3}
 \frac{dx}{\sqrt{(x-r_1)(x-r_2)(r_3-x)}}.
\]
As $t\to-2^+$, $r_1,r_2\to-1$ and $r_3\to2$. Put
$\delta=r_2-r_1$. On a fixed small interval $x=r_2+u$, the third factor
$r_3-x$ is bounded above and below by positive constants, and the integral
is comparable to
\[
 \int_0^{u_0}\frac{du}{\sqrt{u(u+\delta)}}
 =2\operatorname{arsinh}\sqrt{u_0/\delta}\longrightarrow\infty.
\]
Thus it diverges logarithmically, contradicting entire continuation and
proving the lemma.
\end{proof}

\begin{remark}
The singularity at $t=2$ is a singularity of the analytic continuation of a
regular phase-density germ. It is not a claim that the original compactly
supported density is nonzero outside $q(D)=[-11/8,11/8]$.
\end{remark}

\section{Galois disjointness from the actual boundary system}
\subsection{The rank-two homogeneous group}
Eliminating the second coordinate from $Y'=AY$ gives
\[
 y_0''+\frac2z y_0'+\left(\frac5{36z^2}-4\right)y_0=0.
\]
With $w=zy_0$, this is
\[
 w''=\left(4-\frac5{36z^2}\right)w.
\]
Equivalently $y_0=z^{-1/2}v(2z)$, where $v$ satisfies the modified Bessel
equation of order $1/3$. The Bessel differential-Galois classification
states that the group is $\SL_2$ for an order not in
$1/2+\Z$ \cite[Example 1.3.32]{Singer}; an algebraic rank-one twist does
not change this connected derived group. Moreover $\operatorname{tr}A=-2/z$,
so the determinant solution is a constant multiple of $z^{-2}$. Thus the
homogeneous group of our system is exactly $\SL_2$.
Its adjoint formal exponential factors at infinity are
\begin{equation}\label{eq:ell-adjoint}
 1,\ e^{4z},\ e^{-4z}.
\end{equation}

\subsection{The boundary has only one nonabelian simple factor}
The homogeneous rank-two block of the boundary system is the cubic moment
module for $p$. Write
\[
 \alpha=\frac{\sqrt7}{3},\qquad
 \beta=\alpha^3=\frac{7\sqrt7}{27},\qquad c_0=\frac{107}{108}.
\]
Then
\begin{equation}\label{eq:pnormal}
 p(\alpha s+1/3)=c_0+\beta(-s^3+3s).
\end{equation}
The centered cubic homogeneous moments satisfy
\[
 V'=\left[
 \begin{pmatrix}0&2\beta\\2\beta&0\end{pmatrix}
 +\frac1z\begin{pmatrix}-1/3&0\\0&-2/3\end{pmatrix}
 \right]V,
\]
up to the scalar exponential $e^{c_0z}$ and an algebraic constant basis
change. Its first coordinate is again modified Bessel of order $1/3$,
now at $2\beta z$. Its derived group is $\SL_2$, and its adjoint formal
exponential factors are
\begin{equation}\label{eq:boundary-adjoint}
 1,\ e^{4\beta z},\ e^{-4\beta z}.
\end{equation}
All other diagonal factors in the full boundary system are rank one:
the vertical quadratic factor and the endpoint exponentials. The kernel of
the Galois action on these diagonal factors is unipotent, because the
system is block triangular. The reductive quotient therefore has one
simple derived factor $\SL_2$, detected by the cubic block; its remaining
quotient is diagonalizable.

\begin{lemma}[Projective exclusion]\label{lem:projective-exclusion}
The adjoint module of $A$ does not belong to the tensor category of the
boundary system.
\end{lemma}
\begin{proof}
Otherwise that boundary group would have the projective Galois group
$\mathrm{PGL}_2$ of $A$ as a quotient. The unipotent radical maps trivially
to this reductive quotient, so the quotient factors through the reductive
boundary group. Its restriction to the unique simple derived $\SL_2$
is the usual projective quotient up to an automorphism of
$\mathrm{PGL}_2$; such automorphisms over $\C$ are inner. After aligning
this restriction with the cubic block, the two projective maps agree on the whole group, including disconnected
components. Indeed, if $S$ is the connected simple derived factor, both
maps send $gsg^{-1}$ to conjugation of the common image of $s$. Their
values on $g$ therefore differ by an element centralizing
$\mathrm{PGL}_2$. This centralizer is trivial. Thus the two adjoint differential modules
would be isomorphic over $\C(z)$.

Their formal exponential factors \eqref{eq:ell-adjoint} and
\eqref{eq:boundary-adjoint} disagree, since $\beta\ne\pm1$. Such factors
are preserved by differential-module isomorphism. This is a contradiction.
Rank-one endpoint exponentials do not affect an adjoint module, so they
cannot remove the obstruction.
\end{proof}

\begin{theorem}[Complete actual elliptic comparison]\label{thm:elliptic-main}
For the fixed triangle and phase above,
\[
 M_0,M_1\quad\text{are algebraically independent over }K_\partial.
\]
In particular, over the arithmetic boundary field,
\begin{equation}\label{eq:free-algebra}
 K_{\partial,0}[I_D(r\,dx\wedge dy):r\in k[x,y]]
 =K_{\partial,0}[M_0,M_1]\simeq K_{\partial,0}[X_0,X_1].
\end{equation}
The canonical map \eqref{eq:canonical-comparison} is an isomorphism onto the
two-dimensional space of actual polynomial-amplitude periods modulo the
boundary field. For every polynomial amplitude,
\begin{equation}\label{eq:linear-comparison}
 I_D(r\,dx\wedge dy)\in K_\partial
 \quad\Longleftrightarrow\quad
 [r\,dx\wedge dy]=0\text{ in }\Hq.
\end{equation}
All such boundary reductions are the polynomial identities
\eqref{eq:elliptic-normal} integrated by Stokes.
\end{theorem}
\begin{proof}
The homogeneous group is $\SL_2$,
\cref{lem:projective-exclusion} gives projective disjointness, and
\cref{lem:nonzero-quotient} gives the required nonzero actual quotient.
Apply \cref{thm:criterion}. The equality of algebras and the equivalence
follow from the unique normal form and algebraic independence. The
statement over $k$ follows from the stronger statement after adjoining
complex constants.
\end{proof}

\section{Numerical comparison and the correct fiber algebra}
For $\xi\in k^\times$, put
\[
 K_{\partial,\xi}
 =k\bigl(P_0(\xi),P_1(\xi),Q_0(\xi),E_a(\xi),E_b(\xi)\bigr).
\]
It is the field of all polynomial-amplitude values on the three fixed faces
and their endpoints: the one-variable reductions have only powers of $z$
in their denominators.

\begin{theorem}[Algebraic specialization]\label{thm:values}
For every $\xi\in\overline{\Q}^{\times}$, the numbers
$M_0(\xi),M_1(\xi)$ are algebraically independent over
$K_{\partial,\xi}$. Consequently
\[
 K_{\partial,\xi}[I_D(r\,dx\wedge dy)(\xi):r\in k[x,y]]
 \simeq K_{\partial,\xi}[X_0,X_1].
\]
\end{theorem}
\begin{proof}
All coordinates of the boundary system and of its augmentation by
$M_0,M_1$ are $E$-functions, and the combined rational system has no finite
singular point other than zero. Let
\[
 d=\trdeg_{k(z)} k(z)(P_0,P_1,Q_0,E_a,E_b).
\]
By \cref{thm:elliptic-main}, adjoining $M_0,M_1$ raises this degree to
$d+2$. The Siegel--Shidlovskii equality, in the form recalled in
\cite[Theorem 1.1]{Beukers}, applied first to the boundary vector and then
to the combined vector, gives degrees $d$ and $d+2$ at every
$\xi\ne0$. The two additional values therefore have transcendence degree
two over $K_{\partial,\xi}$. The normal-form coefficients are Laurent
polynomials in $z$, so specialization proves the second assertion.
\end{proof}

\begin{corollary}[Value-level twisted exactness]\label{cor:value-exactness}
For $r\in k[x,y]$ and $\xi\in k^\times$, the following are equivalent:
\begin{enumerate}
\item $I_D(r\,dx\wedge dy)(\xi)$ is algebraic over $K_{\partial,\xi}$;
\item it belongs to $K_{\partial,\xi}$;
\item $r\,dx\wedge dy\in(d+\xi\,dq\wedge)\Omega^1_{k[x,y]/k}$.
\end{enumerate}
\end{corollary}
\begin{proof}
Specialize \eqref{eq:elliptic-normal}. A nonzero linear form in two
algebraically independent variables is transcendental over their base
field. Thus the first condition forces $a_r(\xi)=b_r(\xi)=0$, and
specialized polynomial division makes this equivalent to the third
condition. Literal Stokes gives the second condition in that case.
\end{proof}

For clarity, one can also formulate the comparison over a local ring.
Let
\[
 R_\partial=k[z,P_0,P_1,Q_0,E_a,E_b],\quad \pi=z-\xi.
\]
Beukers lifting \cite[Theorems 1.2--1.3]{Beukers} gives
\[
 \ker(\operatorname{ev}_\xi:R_\partial\to\C)=\pi R_\partial.
\]
Indeed, a polynomial expression vanishing at $\xi$ differs from a lifted
functional relation by a multiple of $\pi$. Consequently
$B_\xi=(R_\partial)_{\pi R_\partial}$ is a DVR, and its residue field is
$\kappa_\xi=K_{\partial,\xi}$. The local actual functional algebra is
\[
 B_\xi[M_0,M_1]\simeq B_\xi[X_0,X_1].
\]
Its numerical kernel is exactly $\pi B_\xi[X_0,X_1]$, and the numerical
period map is injective on the \emph{fiber}
\[
 (B_\xi[X_0,X_1])/\pi B_\xi[X_0,X_1]
 =\kappa_\xi[X_0,X_1].
\]
There is no claim of numerical injectivity before quotienting by $\pi$.

\begin{remark}[Formal period operations]
The actual functional and numerical independence theorems above do not use
an unpublished completeness theorem for the boundary. To identify all
relations among boundary symbols with Kontsevich--Zagier operations, use
the separate compact polynomial-line formal comparison of Paper III~\cite{LongPolynomialKZ}.
Then \eqref{eq:elliptic-normal}, ordinary twisted Stokes, and polynomial
freeness prove that every relation for this fixed $D,q$ and arbitrary
polynomial amplitudes is generated by those boundary relations and the
explicit Stokes reductions. There is no remaining determinant generator,
no assumed isotypic normal-function compatibility, and no need to infer
arithmetic constants from punctual support.
\end{remark}


\part{Local formal comparison and its scope}
\section{Boundary specialization without a global normal form}
\label{sec:local-boundary}
The local arithmetic input applies to any finite closed $E$-function
vector; it does not require a global Laurent-polynomial presentation of
its actual function algebra.

\begin{lemma}[Direct boundary specialization]\label{lem:direct-local}
Let $f_1,\ldots,f_N$, with the constant coordinate included when needed,
be a solution of a rational first-order $E$-function system over $k(z)$.
Suppose $\xi\in k^\times$ is ordinary for that system and put
\[
 R=k[z,f_1,\ldots,f_N],\qquad \pi=z-\xi.
\]
Then $\ker(\operatorname{ev}_\xi:R\to\C)=\pi R$.
The localization $B=R_{\pi R}$ is a DVR, its valuation is analytic
vanishing order at $\xi$, and
\[
 \kappa=B/\pi B=\operatorname{Frac}(\operatorname{ev}_\xi R).
\]
\end{lemma}
\begin{proof}
Let $P(z,f(z))\in R$ vanish at $\xi$. Beukers lifting supplies
$Q(z,X)\in k[z,X]$ with $Q(z,f(z))=0$ and
$Q(\xi,X)=P(\xi,X)$; use the constant coordinate to homogenize if
necessary. Thus $P-Q=(z-\xi)H$ in $k[z,X]$ and
$P(z,f(z))=\pi H(z,f(z))$. The reverse containment is immediate.

The ring $R$ is a finitely generated domain, and $\pi R$ is a nonzero
prime ideal. The principal ideal theorem gives it height one. Its
localization is a noetherian local domain with nonzero principal maximal
ideal, hence a DVR \cite[Tag~00PD]{Stacks}. Every local unit has a nonzero value at $\xi$, so
writing an element as $\pi^m$ times a unit identifies the valuation with
analytic order. The residue field is the fraction field of $R/\pi R$,
which is the stated actual value field.
\end{proof}

\begin{remark}[Why no global normal form is used]
The entire line function $h=(e^z-1)/z$ gives the actual algebra
\[
 k[z,h,e^z,e^{-z}]=k[z,h,(1+zh)^{-1}].
\]
Here $z,h$ are algebraically independent: otherwise $e^z=1+zh$ would be
algebraic over $k(z)$, contrary to its essential singularity at infinity.
This ring is normal but is not a polynomial extension of a Laurent
polynomial algebra over $k[z]$. Its nonconstant unit $1+zh$ becomes $1$
modulo $z$, whereas a unit $c\,t^a$ in such a Laurent-polynomial model
can become $1$ modulo $z$ only if it already is $1$. Reduced moment
normal forms over an exponential field do not justify the stronger
unlocalized claim for every actual line-function algebra.
\end{remark}

\section{An abstract local comparison theorem}\label{sec:abstract-local}
Let $B$ and $\kappa$ be as in Lemma~\ref{lem:direct-local}, and let
$K=\operatorname{Frac}B$. Let $U_1,\ldots,U_m$ be actual $E$-functions
which, together with the boundary vector, lie in a finite rational
first-order system ordinary at $\xi$. Assume that an ideal
$J\subset B[X_1,\ldots,X_m]$ is generated by explicitly realized formal
identities and satisfies
\begin{align}
 J K[X]&=\ker(K[X]\longrightarrow K[U]),\label{eq:generic-complete}\\
 B[X]/J&\text{ is }B\text{-flat}.\label{eq:local-flat-hyp}
\end{align}
The phrase ``explicitly realized'' is a separate geometric hypothesis:
for example, the polynomial Stokes certificates of this paper provide
such identities. It is not inferred merely from membership of their
coefficients in the boundary field.

\begin{theorem}[Functional contraction and numerical fiber]
\label{thm:abstract-local}
Under \eqref{eq:generic-complete}--\eqref{eq:local-flat-hyp}, the functional
map $B[X]/J\to\Ocal^{\mathrm{an}}_\xi$ is injective, and so is
\[
 (B[X]/J)\otimes_B\kappa\longrightarrow\C.
\]
Equivalently, the numerical kernel on $B[X]/J$ is
$\pi(B[X]/J)$, not zero before specialization.
\end{theorem}
\begin{proof}
Flatness over the DVR implies torsion-freeness. Thus $B[X]/J$ injects
into its generic fiber, and generic completeness gives
\[
 (JK[X])\cap B[X]=J.
\]
This proves the functional assertion.

For the numerical assertion, clear the finitely many boundary denominators
of a representative $P\in B[X]$. Each denominator has nonzero value at
$\xi$, so this multiplication is by a local unit. The resulting
representative is $\widetilde P(z,f(z),X)$ with
$\widetilde P\in k[z,Y,X]$. If its value at $U(\xi)$ is zero, Beukers
lifting applied to the \emph{combined} closed $E$-function vector gives
$Q(z,Y,X)$ which vanishes functionally and satisfies
\[
 Q(\xi,Y,X)=\widetilde P(\xi,Y,X).
\]
After substituting $Y=f(z)$, the polynomial $Q(z,f(z),X)$ is in the
functional kernel and hence in $J$. The difference
$\widetilde P-Q$ is divisible by $\pi$ in $k[z,Y,X]$. Therefore the
class of $P$ in $B[X]/J$ belongs to $\pi(B[X]/J)$. Division by the
previously cleared unit is legitimate. The reverse containment follows
by evaluation, proving the numerical statement.
\end{proof}

\begin{definition}[Specialized formal algebra]\label{def:specialized-formal}
For a local functional symbol algebra $\mathcal P^{\mathrm{fun,loc}}$
over $B$, its numerical algebra at $\xi$ is
\[
 \mathcal P^{\mathrm{num}}_\xi
 =\mathcal P^{\mathrm{fun,loc}}\otimes_B\kappa
 =\mathcal P^{\mathrm{fun,loc}}/\pi\mathcal P^{\mathrm{fun,loc}}.
\]
Whenever the two hypotheses of Theorem~\ref{thm:abstract-local} are
proved for the stated formal generators, its numerical period map is
injective. The residue field is the boundary-value field, generally not
$k$.
\end{definition}

\section{The abstract multicolumn orbit model}\label{sec:affine-orbits}
The following algebraic calculation is retained for future geometric
applications. Its group hypotheses and its eventual formal realization
are not automatic properties of polynomial flags.

Work over a characteristic-zero differential field $K$ with algebraically
closed constants $k$. For finitely many blocks of ranks $n_\lambda\ge2$,
let $Y_\lambda'=A_\lambda Y_\lambda$ be fundamental matrices and let
\[
 U_\lambda'=A_\lambda U_\lambda+B_\lambda
\]
be matrices of actual particular-solution columns in a common
no-new-constants extension. Suppose the homogeneous group, in these
representations, is a subgroup of $\prod_\lambda\GL_{n_\lambda}$
containing $\prod_\lambda\SL_{n_\lambda}$.

\begin{proposition}[Affine orbit reduction]\label{prop:affine-model}
Constant operations in each column multiplicity space remove the
translation-active columns as independent affine coordinates. Every
remaining column is split, of the form
\[
 U_{\lambda,\mathrm{sp}}=a_\lambda+Y_\lambda C_\lambda,
 \qquad a_\lambda\in\operatorname{Mat}(K),
 \quad C_\lambda\in\operatorname{Mat}(k).
\]
After recording the constant kernel of $C_\lambda$ as affine-linear
relations, let $s_\lambda=\operatorname{rank}C_\lambda\le n_\lambda$.
If $s_\lambda<n_\lambda$, this block has no further polynomial relation.
Only blocks with $s_\lambda=n_\lambda$ contribute determinant coordinates
$\Delta_\lambda=\det(U_\lambda-a_\lambda)$.
\end{proposition}
\begin{proof}
The translation kernel is an invariant vector subgroup of the direct sum
of the standard representations with their column multiplicity spaces.
The product of special-linear groups acts independently, so it has the
form $\bigoplus_\lambda k^{n_\lambda}\otimes R_\lambda$.
Choosing constant bases of the $R_\lambda$ gives independent translations
in the corresponding matrix coordinates. The ideal of the solution
variety is extended from the other coordinates; one does not assert that
an arbitrary multiple of a relation cannot mention a translated variable.

After quotienting the translations, the affine term is a regular
algebraic cocycle of the homogeneous reductive group. Complete
reducibility makes it a coboundary, giving the displayed split form.
Constant column operations record all kernels of $C_\lambda$.
For $s<n$, the $\SL_n$-orbit of a full-column-rank $n\times s$ matrix is
the entire rank-$s$ locus: a $\GL_n$ map can be adjusted to determinant
one on a complementary vector. This locus is dense in the whole matrix
space. For $s=n$, the orbit is exactly the nonzero determinant fiber.
These facts give all assertions.
\end{proof}

Let $T_{\det}\subset\Gm^s$ be the image of the homogeneous group on the
determinants of its square blocks. It may be disconnected. Put
\[
 L=\{a\in\Z^s:x^a|_{T_{\det}}=1\},
 \qquad \eta_a=\Delta^{a^+}/\Delta^{a^-}\in K^\times.
\]

\begin{theorem}[Generic affine relation ideal]\label{thm:generic-orbits}
The complete polynomial relation ideal over $K$ is generated by the
recorded affine-linear relations and
\[
 \Delta^{a^+}-\eta_a\Delta^{a^-},\qquad a\in L.
\]
\end{theorem}
\begin{proof}
The affine reduction leaves full affine factors, square determinant
fibers, and the orbit of the determinant vector under $T_{\det}$.
In Laurent determinant coordinates the latter is a translate of a
closed subgroup of a torus, with partial-character ideal given by $L$.
Its contraction to polynomial coordinates is the lattice ideal generated
by all the displayed binomials, including when $L$ is not saturated;
see \cite{EisenbudSturmfels}.

This also describes its pullback through the determinant maps, without
adding a component at zero determinants. Indeed those maps are flat
(Lemma~\ref{lem:flat-determinant-new}), and the lattice coordinate ring
embeds into its Laurent localization. Tensoring this injection with the
flat matrix coordinate ring embeds the pullback into its localization at
all determinants. There it is exactly the special-linear orbit family;
hence its closure gives the asserted ideal. The Picard--Vessiot torsor
identifies the solution relation ideal with this orbit-closure ideal.
\end{proof}

\subsection{A local model with unit normalization}
Suppose now that the boundary local ring is a DVR $B$ as above, that the
block and column coordinates used in the generic theorem are obtained by
invertible changes over $B$, and that
\[
 a_\lambda\in\operatorname{Mat}(B),\qquad \eta_a\in B^\times.
\]
These are hypotheses on the chosen local model. For an ordinary block
system they follow for an existing split expression from holomorphy of
$U_\lambda,Y_\lambda$ and invertibility of $Y_\lambda(\xi)$: its shift
is $U_\lambda-Y_\lambda C_\lambda$, and the selected full-rank
constant determinant is nonzero. A meromorphic block decomposition must
first be replaced by a $B$-invertible one; this is not inferred from a
generic decomposition alone.

\begin{lemma}[Flat unit-valued lattice quotient]\label{lem:flat-lattice-new}
For a character $\rho:L\to B^\times$, the quotient
\[
 D_\rho=B[\delta_1,\ldots,\delta_s]/
 (\delta^{a^+}-\rho(a)\delta^{a^-}:a\in L)
\]
is free over $B$.
\end{lemma}
\begin{proof}
Declare $u,v\in\N^s$ equivalent if $u-v\in L$. The binomial for
$u-v$, multiplied by their common monomial factor, identifies the two
monomials by the unit $\rho(u-v)$. Choose one representative for every
congruence class. The character identity makes all reductions consistent.
Sending each monomial to the corresponding unit multiple of its chosen
representative defines the inverse to the resulting free-module
presentation. Hence those representatives form a $B$-basis.
\end{proof}

\begin{lemma}[Free determinant morphism]\label{lem:flat-determinant-new}
For every $n\ge1$, $B[x_{ij}]$ is free over $B[\delta]$ under
$\delta\mapsto\det(x_{ij})$.
\end{lemma}
\begin{proof}
It is the quotient of $B[\delta][x_{ij}]$ by
$\det(x_{ij})-\delta$. Choose a monomial order in the $x_{ij}$ for which
the diagonal product is leading. The polynomial is monic. Division by this
single monic polynomial gives a unique remainder in the monomials not
divisible by its leading monomial, which form a $B[\delta]$-basis.
This proof works over an arbitrary commutative coefficient ring.
\end{proof}

\begin{proposition}[Flat normalized affine model]\label{prop:flat-model}
In the specified local coordinates, the quotient by the regular
coordinate-linear relations and unit-valued determinant relations is
$B$-flat and its generic ideal contracts to the same $B$-ideal.
\end{proposition}
\begin{proof}
Translation by the regular shifts is a polynomial automorphism over $B$.
After eliminating the coordinate-linear relations, the quotient is a
polynomial extension of the base change of the product of determinant
maps to $D_\rho$. The two preceding lemmas make it flat over $B$.
Torsion-freeness then proves contraction from the generic fiber.
\end{proof}

\begin{remark}[Admissible clearing of denominators]
Choose $\eta_a=b_a/c_a$ only with $b_a(\xi)c_a(\xi)\ne0$.
Then clearing the denominator multiplies the normalized binomial by a
unit and leaves its local ideal unchanged. With $\pi=z-\xi$, the
presentations $1=1/1=\pi/\pi$ give $\delta-1$ and
$\pi(\delta-1)$, which have different ideals over $B$. In the latter
quotient a nonzero class is killed by $\pi$; its equation disappears
on specialization. Arbitrary nonzero global denominators are therefore
not interchangeable in a local formal presentation.
\end{remark}

\subsection{The remaining geometric hypotheses}
The preceding results prove a general numerical formal theorem only when
the following data have been supplied: an actual polynomial boundary
system; verified homogeneous groups over the actual boundary field;
a generically complete \emph{geometric} affine-linear comparison;
a compatible local lattice; and geometric realizations of any
rank-one determinant identifications. Once these give an explicit ideal
$J$ satisfying \eqref{eq:generic-complete} and
\eqref{eq:local-flat-hyp}, Theorem~\ref{thm:abstract-local} applies.

A determinant character becoming trivial in a differential Galois group
identifies a rank-one differential module. It does not by itself prove
that the identification is generated by pullback, trace, and Stokes on
algebraic chains. One may explicitly enlarge the formal category to admit
those maps, or prove their geometric realization. The two statements
should not be conflated. For the elliptic triangle, there is one actual
rank-two column and the independent-coordinate ideal is zero, so this
additional issue does not arise.
\section{Falsification controls and remaining scope}\label{sec:scope}
\subsection{Punctual support does not determine arithmetic normalization}
For a compactly supported distribution $T$ on the real phase line,
\[
 \mathcal L(T)(z)=\langle T,e^{zt}\rangle,
 \qquad \mathcal L(\partial_tT)=-z\mathcal L(T),
 \qquad \mathcal L(tT)=\partial_z\mathcal L(T).
\]
Products of Laplace transforms correspond to convolution. If $T$ is
supported at $a$, its transform has the form
$\sum_m c_m(-z)^m e^{az}$. Algebraicity of $a$ does not imply
algebraicity of the coefficients $c_m$.

For example, take $q=y^2-x^3+3x$ near its minimum $(-1,0)$ and a small
rational rectangle containing that point. Put $u=x+1$. Then
\[
 q+2=3u^2-u^3+y^2=X^2+Y^2,
 \qquad X=u\sqrt{3-u},\quad Y=y.
\]
The area Jacobian at the origin is $1/\sqrt3$, so the phase density near
$-2$ is $\mathbf1_{t>-2}g(t+2)$ with $g(0)=\pi/\sqrt3$. Its derivative
has the punctual term $(\pi/\sqrt3)\delta_{-2}$, whose Laplace transform
is $(\pi/\sqrt3)e^{-2z}$. This is not an arithmetic $E$-function over
$k$. The actual unsplit compact integral still is an $E$-function; the
nonarithmetic constants need not survive separately in a polynomial
Stokes presentation. The boundary-system theorem keeps that presentation
and makes no inference from support alone.

\subsection{Two elementary comparison controls}
Numerical evaluation on an unspecialized local algebra kills its nonzero
element $z-\xi$. This is why Definition~\ref{def:specialized-formal}
uses the fiber algebra. Likewise, multiplication of an ideal generator
by a nonunit can preserve its generic fiber but destroy its local
flatness, as the determinant countermodel above shows.

Topological relative exactness is not always polynomial relative
exactness. For $q=x^2$, the form $\omega=x\,dy$ is exact on each
component of a regular fiber. If $\omega=dR+a\,dq$ polynomially, its
$dy$ coefficient forces $R=xy+h(x)$; the $dx$ coefficient would require
$y+h'(x)+2xa=0$, impossible in $k[x,y]$. The connected-fiber and
finite-critical-locus hypotheses in the geometric theorem are not
cosmetic.

\subsection{Verification and limits of the current theorem}
The exact certificates in the elliptic proof can be checked independently
by polynomial arithmetic. The repository includes a reproducible script
for monomial normal forms, both actual connection identities, both
Gelfand--Leray identities, the scalar equation, the full boundary system,
and independently integrated Taylor coefficients. Additional rational
inequalities verify the branch and endpoint conditions used in the
analytic continuation. These finite checks are not a kernel-checked proof
of the monodromy or differential-Galois arguments.

The boundary-system construction holds for arbitrary polynomial phases
and compact polynomially parametrized two-chains. The actual independence
criterion has three explicit hypotheses. They are proved here for the
fixed elliptic triangle, with arbitrary polynomial amplitudes and every
nonzero algebraic specialization. The general multiphase and multicolumn
geometric comparison remains subject to the hypotheses isolated in the
local framework. In particular, no theorem for arbitrary isotypic relative
normal functions, arbitrary boundary projective overlap, or unspecified
determinant-line realizations follows from the present proof.

The real-density geometry is not used to impose a real-phase restriction
on the algebraic boundary construction. For genuinely complex phases on
real two-chains the phase support may be two dimensional, and the
one-dimensional density-monodromy argument must be replaced; it is not
silently extended to that setting.

\section*{Acknowledgments and AI-assisted research disclosure}
This manuscript is part of the exponential-integral program developed
with Papers~I--III~\cite{LongEIT,LongLinearKZ,LongPolynomialKZ}. During its
preparation the authors used AI assistants for proof exploration,
adversarial checking, symbolic verification, bibliographic checking, and
language editing. The authors are responsible for the mathematical
statements, proofs, and remaining errors. The September 2026 revision
replaces the earlier general endpoint-completeness claim by the explicit
actual comparison and separately stated general hypotheses above.

\small
\begin{thebibliography}{99}

\bibitem{BakkerMatsushita}
B.~Bakker,
\emph{A short proof of a conjecture of Matsushita},
Adv. Math. \textbf{482} (2025), Paper No.~110554.
\href{https://doi.org/10.1016/j.aim.2025.110554}{doi:10.1016/j.aim.2025.110554}.

\bibitem{Beukers}
F.~Beukers,
\emph{A refined version of the Siegel--Shidlovskii theorem},
Ann.\ of Math. (2) \textbf{163} (2006), no.~1, 369--379.
\href{https://doi.org/10.4007/annals.2006.163.369}{doi:10.4007/annals.2006.163.369}.

\bibitem{BonnetRelativeExactness}
P.~Bonnet,
\emph{Relative exactness modulo a polynomial map and algebraic
$(\mathbb C^p,+)$-actions},
Bull. Soc. Math. France \textbf{131} (2003), no.~3, 373--398.
\href{https://doi.org/10.24033/bsmf.2447}{doi:10.24033/bsmf.2447}.

\bibitem{CommelinHabeggerHuber}
J.~Commelin, P.~Habegger, and A.~Huber,
\emph{Exponential periods and o-minimality},
preprint, arXiv:2007.08280, version~4, 2025.
\href{https://arxiv.org/abs/2007.08280}{arXiv:2007.08280}.

\bibitem{DeligneHodgeIII}
P.~Deligne,
\emph{Th\'eorie de Hodge. III},
Inst. Hautes \'Etudes Sci. Publ. Math. \textbf{44} (1974), 5--77.

\bibitem{DeligneFinitude}
P.~Deligne,
\emph{Un th\'eor\`eme de finitude pour la monodromie},
in \emph{Discrete Groups in Geometry and Analysis},
Progr. Math., vol.~67, Birkh\"auser, Boston, 1987, pp.~1--19.



\bibitem{EisenbudSturmfels}
D.~Eisenbud and B.~Sturmfels,
\emph{Binomial ideals},
Duke Math. J. \textbf{84} (1996), no.~1, 1--45.
\href{https://doi.org/10.1215/S0012-7094-96-08401-X}{doi:10.1215/S0012-7094-96-08401-X}.

\bibitem{DimcaSingularities}
A.~Dimca,
\emph{Singularities and Topology of Hypersurfaces},
Universitext, Springer-Verlag, New York, 1992.

\bibitem{HTT}
R.~Hotta, K.~Takeuchi, and T.~Tanisaki,
\emph{$D$-Modules, Perverse Sheaves, and Representation Theory},
Progr. Math., vol.~236, Birkh\"auser, Boston, 2008.

\bibitem{Kolchin}
E.~R.~Kolchin,
\emph{Differential Algebra and Algebraic Groups},
Pure and Applied Mathematics, vol.~54, Academic Press, New York, 1973.

\bibitem{KontsevichZagier}
M.~Kontsevich and D.~Zagier,
\emph{Periods},
in \emph{Mathematics Unlimited--2001 and Beyond},
Springer, Berlin, 2001, pp.~771--808.
\href{https://doi.org/10.1007/978-3-642-56478-9_39}{doi:10.1007/978-3-642-56478-9\_39}.

\bibitem{LongEIT}
C.~D.~Long and A.-A.~Le~Blanc,
\emph{Algebraic Exponential Integrals: Transcendence and Linear Relations},
companion preprint (Paper~I), August 2026.

\bibitem{LongLinearKZ}
C.~D.~Long and A.-A.~Le~Blanc,
\emph{A Linear Kontsevich--Zagier Theorem for Compact Exponential Periods on
the Affine Line},
companion preprint (Paper~II), August 2026.

\bibitem{LongPolynomialKZ}
C.~D.~Long and A.-A.~Le~Blanc,
\emph{A Polynomial Kontsevich--Zagier Theorem for Compact Exponential Periods
on the Affine Line},
companion preprint (Paper~III), August 2026.


\bibitem{SerreAGCF}
J.-P.~Serre,
\emph{Algebraic Groups and Class Fields},
Graduate Texts in Mathematics, vol.~117, Springer, New York, 1988.

\bibitem{vanDerPutSinger}
M.~van~der~Put and M.~F.~Singer,
\emph{Galois Theory of Linear Differential Equations},
Grundlehren der mathematischen Wissenschaften, vol.~328,
Springer, Berlin, 2003.


\bibitem{Sabbah}
C. Sabbah, \emph{Hypergeometric periods for a tame polynomial},
Portugal. Math. \textbf{63} (2006), 173--226.
Preprint \href{https://arxiv.org/abs/math/9805077}{arXiv:math/9805077},
Proposition 9.1, Remarks 9.4, and Corollary 10.2.
\bibitem{Singer}
M. F. Singer, \emph{Introduction to the Galois theory of linear differential
equations}, \href{https://arxiv.org/abs/0712.4124}{arXiv:0712.4124},
version 2, 2008. In particular Example 1.3.32, Proposition 1.3.33,
and Theorem 1.5.2.

\bibitem{Stacks}
The Stacks Project Authors, \emph{The Stacks Project},
Lemma~10.119.7, Tag~00PD.
\href{https://stacks.math.columbia.edu/tag/00PD}{Tag~00PD}.

\end{thebibliography}
\end{document}
